
\documentclass{article}
\usepackage{colortbl}
\usepackage{makecell}
\usepackage{multirow}
\usepackage{supertabular}

\begin{document}

\newcounter{utterance}

\centering \large Interaction Transcript for game `ta\_mastermind', experiment `Mastermind{-}v0', episode 9 with Dialogue Architects\_\_qwen3.5{-}9b\_\_SCoRe\_Qwen3.5{-}9B.
\vspace{24pt}

{ \footnotesize  \setcounter{utterance}{1}
\setlength{\tabcolsep}{0pt}
\begin{supertabular}{c@{$\;$}|p{.15\linewidth}@{}p{.15\linewidth}p{.15\linewidth}p{.15\linewidth}p{.15\linewidth}p{.15\linewidth}}
    \# & \multicolumn{2}{c}{Player} && \multicolumn{2}{c}{Game Master} \\
    \hline

    \theutterance \stepcounter{utterance}  
    & & & \multicolumn{4}{p{0.6\linewidth}}{
        \cellcolor[rgb]{0.9,0.9,0.9}{
            \makecell[{{p{\linewidth}}}]{
                \texttt{\tiny{[P1$\langle$GM]}}
                \texttt{You are a competitive game player. Make sure you read the game instructions carefully, and always follow the required format.} \\
\\ 
\texttt{{[}GAME{]} You are playing Mastermind.} \\
\texttt{You need to find the code that is 4 digits long, each digit from 1 to 6, with no duplicates.} \\
\texttt{In your response, you can submit your guess in the following format: '{[}2 1 4 5{]}'.} \\
\texttt{After each guess, you will receive feedback in the form of black and white pegs.} \\
\texttt{A black peg indicates a correct digit in the correct position, while a white peg indicates a correct digit in the wrong position.} \\
\texttt{You have 20 turns to guess the code.} \\
\\ 
            }
        }
    }
     \\ \\

    \theutterance \stepcounter{utterance}  
    & \multicolumn{4}{p{0.6\linewidth}}{
        \cellcolor[rgb]{0.561,0.737,0.561}{
            \makecell[{{p{\linewidth}}}]{
                \texttt{\tiny{[P1$\rangle$GM]}}
                \texttt{{[}1 2 3 4{]}} \\
            }
        }
    }
    & & \\ \\

    \theutterance \stepcounter{utterance}  
    & & & \multicolumn{4}{p{0.6\linewidth}}{
        \cellcolor[rgb]{0.9,0.9,0.9}{
            \makecell[{{p{\linewidth}}}]{
                \texttt{\tiny{[P1$\langle$GM]}}
                \texttt{{[}GAME{]} Submitted {[}1 2 3 4{]}. Feedback: 1 black peg(s), 2 white peg(s).} \\
            }
        }
    }
     \\ \\

    \theutterance \stepcounter{utterance}  
    & \multicolumn{4}{p{0.6\linewidth}}{
        \cellcolor[rgb]{0.561,0.737,0.561}{
            \makecell[{{p{\linewidth}}}]{
                \texttt{\tiny{[P1$\rangle$GM]}}
                \texttt{Based on the feedback (1 black, 2 white) from my guess `{[}1 2 3 4{]}`:} \\
\texttt{{-} **1 Black**: One of these digits is in the correct spot.} \\
\texttt{{-} **2 White**: Two other digits from my guess are in the code but in the wrong spots.} \\
\texttt{{-} **Total known digits**: 3 out of 4 are in the code.} \\
\texttt{{-} **Missing digit**: Since the code has no duplicates and uses digits 1{-}6, the one digit I guessed that isn't in the code must be the one that didn't yield a peg. However, I don't know which one yet.} \\
\texttt{{-} **Positional info**: I need to test permutations to isolate the correct digit and its position.} \\
\\ 
\texttt{Let's try a guess that keeps `1` in the first position (assuming it might be the black peg) but swaps the others to test for white pegs, or better yet, let's try to isolate the black peg.} \\
\texttt{Actually, a more efficient strategy is to keep the potential black pegs in place and rotate the others, or swap positions to see if the black peg moves.} \\
\\ 
\texttt{Let's try `{[}1 3 2 5{]}`.} \\
\texttt{{-} Keeps `1` in pos 1.} \\
\texttt{{-} Swaps `2` and `3` (testing if the black peg is 2 or 3, or if the white pegs change).} \\
\texttt{{-} Replaces `4` with `5` (to see if `4` was the non{-}code digit).} \\
\\ 
\texttt{Wait, if `4` was the non{-}code digit, then `1`, `2`, `3` are in the code.} \\
\texttt{If `1` is the black peg, then `2` and `3` are white pegs (in wrong spots).} \\
\texttt{If `2` is the black peg, then `1` and `3` are white pegs.} \\
\texttt{If `3` is the black peg, then `1` and `2` are white pegs.} \\
\\ 
\texttt{Let's try `{[}1 3 2 5{]}`.} \\
\texttt{{-} If `1` is the black peg: `3` and `2` are still in the code (white pegs). `5` is new. Result should be 1B, 2W.} \\
\texttt{{-} If `2` is the black peg: `1` is white, `3` is white. `5` is new. Result should be 1B, 2W.} \\
\texttt{{-} If `3` is the black peg: `1` is white, `2` is white. `5` is new. Result should be 1B, 2W.} \\
\texttt{This doesn't help distinguish the black peg immediately if the set of known digits is the same.} \\
\\ 
\texttt{Let's try to eliminate a digit.} \\
\texttt{Guess: `{[}1 2 5 6{]}`} \\
\texttt{{-} Keeps `1`, `2`.} \\
\texttt{{-} Removes `3`, `4`. Adds `5`, `6`.} \\
\texttt{{-} If `3` was a white peg and `4` was the non{-}code digit: We lose 1 white peg. Result: 1B, 1W.} \\
\texttt{{-} If `4` was a white peg and `3` was the non{-}code digit: We lose 1 white peg. Result: 1B, 1W.} \\
\texttt{{-} If `1` or `2` was the non{-}code digit: We lose a black peg or a white peg.} \\
\\ 
\texttt{Let's go with `{[}1 2 5 6{]}`. This tests if `3` or `4` are in the code.} \\
\\ 
\texttt{{[}1 2 5 6{]}} \\
            }
        }
    }
    & & \\ \\

    \theutterance \stepcounter{utterance}  
    & & & \multicolumn{4}{p{0.6\linewidth}}{
        \cellcolor[rgb]{0.9,0.9,0.9}{
            \makecell[{{p{\linewidth}}}]{
                \texttt{\tiny{[P1$\langle$GM]}}
                \texttt{{[}GAME{]} Submitted {[}1 2 5 6{]}. Feedback: 0 black peg(s), 2 white peg(s).} \\
            }
        }
    }
     \\ \\

    \theutterance \stepcounter{utterance}  
    & \multicolumn{4}{p{0.6\linewidth}}{
        \cellcolor[rgb]{0.561,0.737,0.561}{
            \makecell[{{p{\linewidth}}}]{
                \texttt{\tiny{[P1$\rangle$GM]}}
                \texttt{Based on the feedback from `{[}1 2 5 6{]}` (0 Black, 2 White):} \\
\texttt{{-} **0 Black**: None of the digits are in the correct position. This means `1` is not in position 1, and `2` is not in position 2.} \\
\texttt{{-} **2 White**: Two digits from my guess are in the code but in the wrong positions.} \\
\texttt{{-} **Comparison with previous guess `{[}1 2 3 4{]}` (1B, 2W)**:} \\
\texttt{    {-} Previous guess had `1, 2, 3, 4`. Current has `1, 2, 5, 6`.} \\
\texttt{    {-} The common digits are `1` and `2`.} \\
\texttt{    {-} If `1` and `2` were the two white pegs in the current guess, then `3` and `4` must be the non{-}code digits.} \\
\texttt{    {-} Let's verify: If `3` and `4` are out, then the code contains `1, 2, 5, 6`.} \\
\texttt{    {-} In the first guess `{[}1 2 3 4{]}`, if the code is a permutation of `1, 2, 5, 6`, then `3` and `4` contribute 0 pegs. That leaves `1` and `2` to account for the 1 Black and 2 White pegs? No, that's impossible. 1 Black + 2 White = 3 pegs total. If only `1` and `2` are in the code, we can get at most 2 pegs (1B+1W or 0B+2W).} \\
\texttt{    {-} Therefore, my assumption that `3` and `4` are the non{-}code digits is **incorrect**.} \\
\texttt{    {-} One of `3` or `4` must be in the code, and one of `5` or `6` must be in the code.} \\
\texttt{    {-} Let's re{-}evaluate the overlap.} \\
\texttt{    {-} Guess 1: `1, 2, 3, 4` {-}> 1B, 2W. (3 correct digits).} \\
\texttt{    {-} Guess 2: `1, 2, 5, 6` {-}> 0B, 2W. (2 correct digits).} \\
\texttt{    {-} The set of correct digits in Guess 1 is \{A, B, C\}. The set in Guess 2 is \{D, E\}.} \\
\texttt{    {-} Intersection of sets: \{1, 2\}.} \\
\texttt{    {-} If the intersection \{1, 2\} contained 2 correct digits, Guess 2 would have at least 2 pegs (which it does). But Guess 1 would have at most 2 pegs from \{1, 2\} plus whatever `3` or `4` contributed.} \\
\texttt{    {-} Wait, if \{1, 2\} are the only common correct digits, then Guess 1 has 2 correct digits. But Guess 1 gave 3 pegs (1B+2W). This implies there is a 3rd correct digit in Guess 1 that is NOT in Guess 2. That digit must be `3` or `4`.} \\
\texttt{    {-} So, the code contains two digits from \{1, 2\} and one digit from \{3, 4\}.} \\
\texttt{    {-} Also, Guess 2 has 2 correct digits. Since it shares \{1, 2\} with Guess 1, and we know one of \{3, 4\} is in the code, the code must contain exactly one of \{1, 2\} and one of \{5, 6\}?} \\
\texttt{    {-} Let's try a different angle.} \\
\texttt{    {-} Let \$S\_1 = \textbackslash{}\{1, 2, 3, 4\textbackslash{}\}\$, Pegs = 3.} \\
\texttt{    {-} Let \$S\_2 = \textbackslash{}\{1, 2, 5, 6\textbackslash{}\}\$, Pegs = 2.} \\
\texttt{    {-} Difference \$S\_1 \textbackslash{}setminus S\_2 = \textbackslash{}\{3, 4\textbackslash{}\}\$. Difference \$S\_2 \textbackslash{}setminus S\_1 = \textbackslash{}\{5, 6\textbackslash{}\}\$.} \\
\texttt{    {-} If we remove `3` and `4` from \$S\_1\$, we are left with \$\textbackslash{}\{1, 2\textbackslash{}\}\$. If \$\textbackslash{}\{1, 2\textbackslash{}\}\$ were the only correct digits, max pegs = 2. But we got 3. So at least one of `3` or `4` is in the code.} \\
\texttt{    {-} If we remove `1` and `2` from \$S\_2\$, we are left with \$\textbackslash{}\{5, 6\textbackslash{}\}\$. If \$\textbackslash{}\{5, 6\textbackslash{}\}\$ were the only correct digits, max pegs = 2. We got 2. This is consistent.} \\
\texttt{    {-} So, the code contains exactly one digit from \$\textbackslash{}\{3, 4\textbackslash{}\}\$ and exactly two digits from \$\textbackslash{}\{1, 2, 5, 6\textbackslash{}\}\$.} \\
\texttt{    {-} Actually, let's look at the counts.} \\
\texttt{    {-} Total correct digits = 4.} \\
\texttt{    {-} From \$S\_1\$: 3 correct.} \\
\texttt{    {-} From \$S\_2\$: 2 correct.} \\
\texttt{    {-} Overlap \$\textbackslash{}\{1, 2\textbackslash{}\}\$.} \\
\texttt{    {-} Let \$k\$ be the number of correct digits in \$\textbackslash{}\{1, 2\textbackslash{}\}\$.} \\
\texttt{    {-} Let \$m\$ be the number of correct digits in \$\textbackslash{}\{3, 4\textbackslash{}\}\$.} \\
\texttt{    {-} Let \$n\$ be the number of correct digits in \$\textbackslash{}\{5, 6\textbackslash{}\}\$.} \\
\texttt{    {-} \$k + m = 3\$ (from Guess 1).} \\
\texttt{    {-} \$k + n = 2\$ (from Guess 2).} \\
\texttt{    {-} Possible values for \$k\$:} \\
\texttt{        {-} If \$k=2\$: Then \$m=1\$, \$n=0\$. Code has both 1 and 2, one of 3/4, none of 5/6.} \\
\texttt{        {-} If \$k=1\$: Then \$m=2\$ (impossible, only 2 digits), \$n=1\$. Code has one of 1/2, both 3/4, one of 5/6.} \\
\texttt{        {-} If \$k=0\$: Then \$m=3\$ (impossible).} \\
\texttt{    {-} So either (\$k=2, m=1, n=0\$) or (\$k=1, m=2, n=1\$).} \\
\texttt{    {-} Case A: \$k=2\$. Digits 1 and 2 are in the code. One of 3 or 4 is in. None of 5 or 6 are in.} \\
\texttt{        {-} If this is true, then in Guess 2 (`1 2 5 6`), the correct digits are 1 and 2. That's 2 digits. Pegs = 2. This matches.} \\
\texttt{        {-} In Guess 1 (`1 2 3 4`), correct digits are 1, 2, and one of 3/4. That's 3 digits. Pegs = 3. This matches.} \\
\texttt{    {-} Case B: \$k=1\$. One of 1 or 2 is in. Both 3 and 4 are in. One of 5 or 6 is in.} \\
\texttt{        {-} If this is true, in Guess 2 (`1 2 5 6`), correct digits are one of 1/2 and one of 5/6. Total 2. Pegs = 2. Matches.} \\
\texttt{        {-} In Guess 1 (`1 2 3 4`), correct digits are one of 1/2 and both 3/4. Total 3. Pegs = 3. Matches.} \\
\texttt{    {-} We need to distinguish between Case A and Case B.} \\
\texttt{    {-} Let's test a guess that includes `3` and `4` to see if we get 2 white pegs (supporting Case B) or 1 white peg (supporting Case A).} \\
\texttt{    {-} But wait, we also have positional info.} \\
\texttt{    {-} Guess 1: `1 2 3 4` {-}> 1B, 2W.} \\
\texttt{    {-} Guess 2: `1 2 5 6` {-}> 0B, 2W.} \\
\texttt{    {-} In Guess 2, `1` is not in pos 1, `2` is not in pos 2.} \\
\texttt{    {-} In Guess 1, one of them IS in the correct position.} \\
\texttt{    {-} If Case A is true (1 and 2 are in code):} \\
\texttt{        {-} In Guess 1, since 1 and 2 are in code, and we have 1B, 2W.} \\
\texttt{        {-} If `1` is the black peg (pos 1), then `2` and the other of \{3,4\} are white.} \\
\texttt{        {-} If `2` is the black peg (pos 2), then `1` and the other of \{3,4\} are white.} \\
\texttt{        {-} In Guess 2 (`1 2 5 6`), `1` is not pos 1, `2` is not pos 2. So `1` and `2` must be white pegs. Total 2W. This fits perfectly.} \\
\texttt{    {-} If Case B is true (one of 1/2, both 3/4, one of 5/6):} \\
\texttt{        {-} In Guess 2, we have 2 correct digits (one from 1/2, one from 5/6). Both are white (0B).} \\
\texttt{        {-} In Guess 1, we have 3 correct digits (one from 1/2, both 3/4).} \\
\texttt{        {-} We need 1B, 2W.} \\
\texttt{        {-} If the correct digit from \{1,2\} is in the right spot in Guess 1, we get 1B. The other two (one from 1/2, both 3/4) must be white. But there are 3 correct digits total. 1B + 2W = 3. This works too.} \\
\texttt{    {-} So both cases are still possible based on counts alone.} \\
\texttt{    {-} Let's try to force a specific configuration.} \\
\texttt{    {-} Let's guess `{[}3 4 1 2{]}`.} \\
\texttt{    {-} If Case A (1, 2, X, Y where X,Y in \{3,4\} is false, actually 1,2, one of 3/4):} \\
\texttt{        {-} Digits: 1, 2, 3, 4.} \\
\texttt{        {-} Positions: 3, 4, 1, 2.} \\
\texttt{        {-} If code is `1 3 2 4` (example):} \\
\texttt{            {-} Guess 1 `1 2 3 4`: 1(B), 2(W), 3(W), 4(W)? No, 4 is in code? No, code has 1,2,3,4? No, code has 1,2,3.} \\
\texttt{            {-} Let's assume code is `1 3 2 5` (Case A: 1,2,3 in code, 5 not? No, Case A said 5,6 not in code).} \\
\texttt{            {-} If Case A: Code has 1, 2, and one of 3/4. Say `1 3 2 4` is impossible (duplicates? No, 1,2,3,4 are distinct). But we need 4 digits. If 5,6 are out, code must be permutation of 1,2,3,4.} \\
\texttt{            {-} If code is permutation of 1,2,3,4:} \\
\texttt{                {-} Guess 1 `1 2 3 4`: 1B, 2W.} \\
\texttt{                {-} Guess 2 `1 2 5 6`: 0B, 2W. (Matches 1,2 being white).} \\
\texttt{                {-} So if code is a permutation of 1,2,3,4, then 5 and 6 are definitely out.} \\
\texttt{            {-} If Case B: Code has one of 1/2, both 3/4, one of 5/6.} \\
\texttt{                {-} Example: `3 4 1 5`.} \\
\texttt{                {-} Guess 1 `1 2 3 4`: 1(W), 2(out), 3(B?), 4(B?).} \\
\texttt{                    {-} If code `3 4 1 5`:} \\
\texttt{                    {-} Guess 1 `1 2 3 4`: 1 is at pos 3 (W), 2 is out, 3 is at pos 1 (W), 4 is at pos 2 (W). Total 3W. But we got 1B, 2W.} \\
\texttt{                    {-} So `3 4 1 5` doesn't work for Guess 1.} \\
\texttt{                {-} Try `3 1 4 2` (Case B: 3,4,1,2? No, Case B says one of 1/2. Here we have 1 and 2. So Case B fails if we have both 1 and 2).} \\
\texttt{                {-} Case B requires exactly one of 1 or 2.} \\
\texttt{                {-} Try code `3 4 5 1`. (Has 3,4,5,1. One of 1/2 is 1. Both 3/4. One of 5/6 is 5).} \\
\texttt{                {-} Guess 1 `1 2 3 4`:} \\
\texttt{                    {-} 1 is at pos 4 (W).} \\
\texttt{                    {-} 2 is out.} \\
\texttt{                    {-} 3 is at pos 1 (W).} \\
\texttt{                    {-} 4 is at pos 2 (W).} \\
\texttt{                    {-} Result: 0B, 3W. We need 1B, 2W. Fail.} \\
\texttt{                {-} Try code `3 4 5 2`.} \\
\texttt{                    {-} Guess 1 `1 2 3 4`:} \\
\texttt{                    {-} 1 out.} \\
\texttt{                    {-} 2 at pos 2 (B).} \\
\texttt{                    {-} 3 at pos 1 (W).} \\
\texttt{                    {-} 4 at pos 2? No, 4 is at pos 2 in guess, code has 4 at pos 2? No, code `3 4 5 2`. 4 is at pos 2. Guess 4 is at pos 4. So 4 is W.} \\
\texttt{                    {-} Result: 1B (2), 2W (3, 4). Matches!} \\
\texttt{                    {-} Check Guess 2 `1 2 5 6` against `3 4 5 2`:} \\
\texttt{                    {-} 1 out.} \\
\texttt{                    {-} 2 at pos 2 (B). But we need 0B. Fail.} \\
\texttt{    {-} It seems Case B is very restrictive and likely impossible given the 0B in Guess 2.} \\
\texttt{    {-} If Guess 2 has 0B, then `1` is not at pos 1, `2` is not at pos 2, `5` is not at pos 3, `6` is not at pos 4.} \\
\texttt{    {-} If Case B holds (one of 1/2, both 3/4, one of 5/6):} \\
\texttt{        {-} We need 1B in Guess 1.} \\
\texttt{        {-} If the black peg in Guess 1 is `3` (pos 3) or `4` (pos 4).} \\
\texttt{        {-} If `3` is B (pos 3): Code has 3 at pos 3.} \\
\texttt{        {-} If `4` is B (pos 4): Code has 4 at pos 4.} \\
\texttt{        {-} Let's assume Case A is more likely: The code is a permutation of 1, 2, 3, 4.} \\
\texttt{        {-} Why? Because if 5 or 6 were in, we'd have to explain the counts carefully.} \\
\texttt{        {-} Let's try to confirm if 5 and 6 are out.} \\
\texttt{        {-} Guess `{[}3 4 1 2{]}`.} \\
\texttt{        {-} If code is permutation of 1,2,3,4:} \\
\texttt{            {-} Digits are 1,2,3,4.} \\
\texttt{            {-} Guess 1 `1 2 3 4` {-}> 1B, 2W.} \\
\texttt{            {-} Guess 2 `1 2 5 6` {-}> 0B, 2W (1,2 are white).} \\
\texttt{            {-} Guess 3 `3 4 1 2`:} \\
\texttt{                {-} If code is `1 3 2 4`:} \\
\texttt{                    {-} 3 at pos 1 (W), 4 at pos 2 (W), 1 at pos 3 (W), 2 at pos 4 (W). 4W. No.} \\
\texttt{                {-} If code is `2 1 4 3`:} \\
\texttt{                    {-} 3 at pos 1 (W), 4 at pos 2 (W), 1 at pos 3 (W), 2 at pos 4 (W). 4W. No.} \\
\texttt{                {-} If code is `1 2 3 4` (impossible, 1B 2W).} \\
\texttt{                {-} If code is `2 3 4 1`:} \\
\texttt{                    {-} 3 at pos 1 (W), 4 at pos 2 (W), 1 at pos 3 (W), 2 at pos 4 (W). 4W.} \\
\texttt{                {-} If code is `3 1 2 4`:} \\
\texttt{                    {-} 3 at pos 1 (B). 4 at pos 2 (W). 1 at pos 3 (W). 2 at pos 4 (W). 1B 3W. No.} \\
\texttt{                {-} If code is `3 4 1 2`:} \\
\texttt{                    {-} 3 at pos 1 (B). 4 at pos 2 (B). 1 at pos 3 (W). 2 at pos 4 (W). 2B 2W. No.} \\
\texttt{                {-} If code is `4 3 2 1`:} \\
\texttt{                    {-} 3 at pos 1 (W), 4 at pos 2 (W), 1 at pos 3 (W), 2 at pos 4 (W). 4W.} \\
\texttt{        {-} This suggests my manual simulation of permutations is getting messy. Let's use a logical deduction.} \\
\texttt{        {-} We know 1 and 2 are in the code (from Case A being the only one that fits the "0B in Guess 2" constraint easily? Or maybe not).} \\
\texttt{        {-} Let's try a guess that isolates the black peg.} \\
\texttt{        {-} Guess `{[}3 1 2 4{]}`.} \\
\texttt{        {-} If 1 is the black peg (pos 1):} \\
\texttt{            {-} Guess 1 `1 2 3 4`: 1 is B. 2,3,4 are W? No, 1B 2W. So 2 of \{2,3,4\} are W.} \\
\texttt{            {-} Guess 2 `1 2 5 6`: 1 is W (pos 1, but code has 1 elsewhere? No, if 1 is B in Guess 1, it's at pos 1. In Guess 2, 1 is at pos 1, so it would be B. But Guess 2 is 0B. Contradiction).} \\
\texttt{            {-} **CRITICAL DEDUCTION**:} \\
\texttt{                {-} In Guess 1 (`1 2 3 4`), we have 1 Black.} \\
\texttt{                {-} In Guess 2 (`1 2 5 6`), we have 0 Black.} \\
\texttt{                {-} The digits `1` and `2` are in the same positions in both guesses.} \\
\texttt{                {-} If `1` was the black peg in Guess 1, it would be at pos 1. Then in Guess 2, `1` is at pos 1, so it should be a black peg. But Guess 2 has 0 Black.} \\
\texttt{                {-} Therefore, `1` is NOT the black peg.} \\
\texttt{                {-} Similarly, if `2` was the black peg in Guess 1, it would be at pos 2. Then in Guess 2, `2` is at pos 2, so it should be a black peg. But Guess 2 has 0 Black.} \\
\texttt{                {-} Therefore, `2` is NOT the black peg.} \\
\texttt{                {-} So the black peg in Guess 1 must be `3` or `4`.} \\
\texttt{                {-} And `1` and `2` must be white pegs in Guess 1 (since they are not the black peg, and they are in the code? Wait, are they in the code?).} \\
\texttt{                {-} If `1` and `2` are not in the code, then Guess 1 has 0 white pegs from them. Then we need 2 white pegs from `3` and `4`. Impossible (only 2 digits).} \\
\texttt{                {-} So `1` and `2` MUST be in the code.} \\
\texttt{                {-} Since `1` and `2` are in the code but not the black peg, they must be white pegs in Guess 1.} \\
\texttt{                {-} So in Guess 1: `1` is White, `2` is White. `3` or `4` is Black. The other of `3`/`4` is not in the code (since we only have 3 pegs total: 1B + 2W).} \\
\texttt{                {-} So the code contains `1`, `2`, and one of `3` or `4`.} \\
\texttt{                {-} Now look at Guess 2 (`1 2 5 6`): 0 Black, 2 White.} \\
\texttt{                {-} `1` and `2` are in the code. In Guess 2, they are at pos 1 and 2.} \\
\texttt{                {-} We know `1` is not at pos 1 (from Guess 1 logic? No, we know `1` is not the black peg in Guess 1, which means `1` is not at pos 1 in the code. Correct).} \\
\texttt{                {-} We know `2` is not at pos 2 in the code.} \\
\texttt{                {-} So in Guess 2, `1` is White (correct digit, wrong pos) and `2` is White.} \\
\texttt{                {-} This accounts for the 2 White pegs in Guess 2.} \\
\texttt{                {-} This implies `5` and `6` are NOT in the code.} \\
\texttt{                {-} So the code consists of `1`, `2`, and one of `3` or `4`.} \\
\texttt{                {-} We need to find which one (`3` or `4`) is in the code, and where `1` and `2` are.} \\
\texttt{                {-} From Guess 1: `1` is W, `2` is W. One of `3`/`4` is B.} \\
\texttt{                {-} If `3` is B (pos 3): Code has `3` at pos 3. `1` and `2` are in \{pos 1, 2, 4\} but not} \\
            }
        }
    }
    & & \\ \\

    \theutterance \stepcounter{utterance}  
    & & & \multicolumn{4}{p{0.6\linewidth}}{
        \cellcolor[rgb]{0.9,0.9,0.9}{
            \makecell[{{p{\linewidth}}}]{
                \texttt{\tiny{[P1$\langle$GM]}}
                \texttt{{[}GAME{]} Submitted {[}3 1 2 4{]}. Feedback: 0 black peg(s), 3 white peg(s).} \\
            }
        }
    }
     \\ \\

    \theutterance \stepcounter{utterance}  
    & \multicolumn{4}{p{0.6\linewidth}}{
        \cellcolor[rgb]{0.561,0.737,0.561}{
            \makecell[{{p{\linewidth}}}]{
                \texttt{\tiny{[P1$\rangle$GM]}}
                \texttt{Based on the feedback from `{[}3 1 2 4{]}` (0 Black, 3 White):} \\
\texttt{{-} **0 Black**: None of the digits are in the correct position.} \\
\texttt{    {-} `3` is not at pos 1.} \\
\texttt{    {-} `1` is not at pos 2.} \\
\texttt{    {-} `2` is not at pos 3.} \\
\texttt{    {-} `4` is not at pos 4.} \\
\texttt{{-} **3 White**: Three digits from my guess are in the code but in the wrong positions.} \\
\texttt{{-} **Deduction**:} \\
\texttt{    {-} My previous guess `{[}1 2 3 4{]}` gave 1 Black, 2 White.} \\
\texttt{    {-} My guess `{[}3 1 2 4{]}` gives 0 Black, 3 White.} \\
\texttt{    {-} The set of digits in both guesses is `\{1, 2, 3, 4\}`.} \\
\texttt{    {-} Since the number of white pegs increased from 2 to 3, and the black peg count dropped to 0, it implies that the "Black" peg in the first guess was actually a digit that moved to a position where it is now counted as a White peg (or the configuration changed such that no digit is in the correct spot, but more digits are present).} \\
\texttt{    {-} Wait, if the set of digits is identical, the number of correct digits (pegs) should be constant unless the definition of "correct" changes. But the set of correct digits is fixed.} \\
\texttt{    {-} Let's re{-}evaluate the "Black" deduction.} \\
\texttt{    {-} In `{[}1 2 3 4{]}`: 1B, 2W. Total 3 correct digits.} \\
\texttt{    {-} In `{[}3 1 2 4{]}`: 0B, 3W. Total 3 correct digits.} \\
\texttt{    {-} This confirms the code contains exactly 3 digits from `\{1, 2, 3, 4\}` and 1 digit from outside (which we previously deduced must be one of 3 or 4? No, let's re{-}verify the "outside" digit).} \\
\texttt{    {-} Let's look at the intersection with `{[}1 2 5 6{]}` (0B, 2W).} \\
\texttt{        {-} `{[}1 2 5 6{]}` has 2 correct digits.} \\
\texttt{        {-} `{[}1 2 3 4{]}` has 3 correct digits.} \\
\texttt{        {-} Overlap `\{1, 2\}`.} \\
\texttt{        {-} If `\{1, 2\}` are both correct, then `{[}1 2 5 6{]}` has 2 correct digits (1 and 2). This matches the 2W feedback (since 0B, they must be wrong positions).} \\
\texttt{        {-} If `\{1, 2\}` are both correct, then `{[}1 2 3 4{]}` has at least 2 correct digits. It has 3. So one of `\{3, 4\}` is correct.} \\
\texttt{        {-} So the code contains `1`, `2`, and one of `\{3, 4\}`. The fourth digit is NOT in `\{1, 2, 3, 4\}`.} \\
\texttt{        {-} Wait, if the code has 4 digits, and 3 are from `\{1, 2, 3, 4\}`, the 4th is from `\{5, 6\}`.} \\
\texttt{        {-} But in `{[}1 2 5 6{]}`, we got 2W. If the code has `1`, `2`, and say `3`, then `{[}1 2 5 6{]}` only has `1` and `2` in common. That's 2 digits. Matches 2W.} \\
\texttt{        {-} So the code is a permutation of `\{1, 2, X, Y\}` where `X` is in `\{3, 4\}` and `Y` is in `\{5, 6\}`?} \\
\texttt{        {-} Let's check `{[}3 1 2 4{]}` again. Digits: `3, 1, 2, 4`.} \\
\texttt{        {-} If the code has `1, 2, 3, 4` (i.e., Y is in 3/4), then `{[}3 1 2 4{]}` has 4 correct digits. But we got 3W. This implies one digit is NOT in the code.} \\
\texttt{        {-} So the code has 3 digits from `\{1, 2, 3, 4\}` and 1 digit from `\{5, 6\}`.} \\
\texttt{        {-} Which one is missing from `\{1, 2, 3, 4\}`?} \\
\texttt{        {-} Compare `{[}1 2 3 4{]}` (3 pegs) and `{[}3 1 2 4{]}` (3 pegs).} \\
\texttt{        {-} Both have 3 correct digits.} \\
\texttt{        {-} In `{[}1 2 3 4{]}`, positions are 1, 2, 3, 4.} \\
\texttt{        {-} In `{[}3 1 2 4{]}`, positions are 3, 1, 2, 4.} \\
\texttt{        {-} The digit `4` is at pos 4 in both.} \\
\texttt{        {-} If `4` were in the code, and it was the one missing from the set in the other guess... no, the set is the same.} \\
\texttt{        {-} Let's assume the missing digit is `4`. Then the code has `1, 2, 3` and one of `5, 6`.} \\
\texttt{            {-} Guess `{[}1 2 3 4{]}`: Contains `1, 2, 3`. 3 correct digits. Feedback 1B, 2W.} \\
\texttt{            {-} Guess `{[}3 1 2 4{]}`: Contains `1, 2, 3`. 3 correct digits. Feedback 0B, 3W.} \\
\texttt{            {-} This is consistent.} \\
\texttt{        {-} Let's assume the missing digit is `3`. Then the code has `1, 2, 4` and one of `5, 6`.} \\
\texttt{            {-} Guess `{[}1 2 3 4{]}`: Contains `1, 2, 4`. 3 correct digits. Feedback 1B, 2W.} \\
\texttt{            {-} Guess `{[}3 1 2 4{]}`: Contains `1, 2, 4`. 3 correct digits. Feedback 0B, 3W.} \\
\texttt{            {-} This is also consistent.} \\
\texttt{        {-} Let's assume the missing digit is `1` or `2`.} \\
\texttt{            {-} If missing `1`: `{[}1 2 3 4{]}` has `2, 3, 4` (3 digits). `{[}3 1 2 4{]}` has `2, 3, 4` (3 digits). Consistent.} \\
\texttt{            {-} If missing `2`: Same.} \\
\texttt{        {-} So we don't know which digit is missing yet.} \\
\texttt{        {-} However, we know `1` and `2` are in the code because in `{[}1 2 5 6{]}` we got 2W. If `1` or `2` were missing, we would have at most 1 correct digit from `\{1, 2\}` plus maybe `5` or `6`.} \\
\texttt{            {-} If code has `5, 6` and one of `3, 4`: `{[}1 2 5 6{]}` has `5, 6` (2 digits). 2W. Matches.} \\
\texttt{            {-} But then `{[}1 2 3 4{]}` would have only one of `3, 4`. Total 1 correct digit. But we got 3 pegs. Contradiction.} \\
\texttt{            {-} Therefore, `1` and `2` MUST be in the code.} \\
\texttt{        {-} So the code contains `1`, `2`, and one of `\{3, 4\}`, and one of `\{5, 6\}`.} \\
\texttt{        {-} Now, look at `{[}3 1 2 4{]}` (0B, 3W).} \\
\texttt{            {-} Digits present: `1, 2, 3, 4`.} \\
\texttt{            {-} If the code has `1, 2, 3, 4` (i.e. the 4th digit is from 3/4, meaning the set is \{1,2,3,4\}), then we have 4 correct digits. But we got 3W. So the code does NOT contain all of \{1,2,3,4\}.} \\
\texttt{            {-} So exactly one of `\{1, 2, 3, 4\}` is NOT in the code.} \\
\texttt{            {-} We know `1` and `2` are in. So either `3` or `4` is the missing one.} \\
\texttt{            {-} Case A: Missing `3`. Code has `1, 2, 4, X` (X in 5/6).} \\
\texttt{                {-} Guess `{[}1 2 3 4{]}`: Has `1, 2, 4`. (3 correct).} \\
\texttt{                {-} Guess `{[}3 1 2 4{]}`: Has `1, 2, 4`. (3 correct).} \\
\texttt{            {-} Case B: Missing `4`. Code has `1, 2, 3, X` (X in 5/6).} \\
\texttt{                {-} Guess `{[}1 2 3 4{]}`: Has `1, 2, 3`. (3 correct).} \\
\texttt{                {-} Guess `{[}3 1 2 4{]}`: Has `1, 2, 3`. (3 correct).} \\
\texttt{        {-} We need to distinguish between Case A and Case B.} \\
\texttt{        {-} Let's look at the positions in `{[}3 1 2 4{]}` (0B).} \\
\texttt{            {-} `3` is at pos 1. If `3` is in the code, it's not at pos 1.} \\
\texttt{            {-} `1` is at pos 2. If `1` is in the code, it's not at pos 2.} \\
\texttt{            {-} `2` is at pos 3. If `2` is in the code, it's not at pos 3.} \\
\texttt{            {-} `4` is at pos 4. If `4` is in the code, it's not at pos 4.} \\
\texttt{        {-} Let's look at `{[}1 2 3 4{]}` (1B, 2W).} \\
\texttt{            {-} One of `1, 2, 3, 4` is in the correct position.} \\
\texttt{            {-} We know `1` is not at pos 1 (from `{[}3 1 2 4{]}`? No, in `{[}3 1 2 4{]}`, `1` is at pos 2. If `1` was at pos 1 in code, it would be W in `{[}3 1 2 4{]}` (since pos 1 has 3). Wait.} \\
\texttt{            {-} In `{[}3 1 2 4{]}`:} \\
\texttt{                {-} Pos 1 has `3`. If code has `1` at pos 1, then `3` is W (if 3 in code) or out.} \\
\texttt{                {-} Pos 2 has `1`. If code has `1` at pos 1, then `1` at pos 2 is W.} \\
\texttt{            {-} Let's test the hypothesis: **Is `3` in the code?**} \\
\texttt{                {-} If `3` is NOT in the code (Case A):} \\
\texttt{                    {-} Code has `1, 2, 4, X`.} \\
\texttt{                    {-} In `{[}1 2 3 4{]}`: `1, 2, 4` are correct. `3` is out.} \\
\texttt{                    {-} In `{[}3 1 2 4{]}`: `1, 2, 4` are correct. `3` is out.} \\
\texttt{                    {-} In `{[}1 2 3 4{]}`, we have 1B, 2W.} \\
\texttt{                        {-} Possible positions for `1, 2, 4`:} \\
\texttt{                        {-} `1` cannot be pos 1 (from `{[}3 1 2 4{]}`? No, in `{[}3 1 2 4{]}`, `1` is at pos 2. If `1` was at pos 1, it's W. If `1` was at pos 2, it would be B. But we got 0B. So `1` is NOT at pos 2.} \\
\texttt{                        {-} `2` cannot be pos 3 (from `{[}3 1 2 4{]}`, `2` is at pos 3, 0B {-}> `2` not at pos 3).} \\
\texttt{                        {-} `4` cannot be pos 4 (from `{[}3 1 2 4{]}`, `4` is at pos 4, 0B {-}> `4` not at pos 4).} \\
\texttt{                        {-} So in `{[}1 2 3 4{]}`:} \\
\texttt{                            {-} `1` is not at pos 1? We don't know yet.} \\
\texttt{                            {-} `2` is not at pos 2? We don't know yet.} \\
\texttt{                            {-} `4` is not at pos 4.} \\
\texttt{                            {-} We need 1B.} \\
\texttt{                            {-} If `1` is B (pos 1): Then `2` and `4` are W.} \\
\texttt{                                {-} Check `{[}3 1 2 4{]}`: `1` at pos 2 (W, ok). `2` at pos 3 (W, ok). `4` at pos 4 (W, ok). `3` out. Total 3W. Matches.} \\
\texttt{                            {-} If `2` is B (pos 2): Then `1` and `4` are W.} \\
\texttt{                                {-} Check `{[}3 1 2 4{]}`: `1` at pos 2 (B? No, code has `2` at pos 2. Guess has `1`. So `1` is W or out. If `1` is in code, it's W. `2` is at pos 3 in guess. Code has `2` at pos 2. So `2` is W. `4` at pos 4 (W). `3` out. Total 3W. Matches.} \\
\texttt{                            {-} If `4` is B (pos 3? No, `4` is at pos 4 in guess 1. If `4` is B, it's at pos 4. But we know `4` is not at pos 4 from `{[}3 1 2 4{]}`. So `4` cannot be B in guess 1).} \\
\texttt{                            {-} So if Case A (`3` out), `1` or `2` is the black peg.} \\
\texttt{                {-} If `4` is NOT in the code (Case B):} \\
\texttt{                    {-} Code has `1, 2, 3, X`.} \\
\texttt{                    {-} In `{[}1 2 3 4{]}`: `1, 2, 3` are correct. `4` is out.} \\
\texttt{                    {-} In `{[}3 1 2 4{]}`: `1, 2, 3` are correct. `4` is out.} \\
\texttt{                    {-} Constraints from `{[}3 1 2 4{]}` (0B):} \\
\texttt{                        {-} `3` not at pos 1.} \\
\texttt{                        {-} `1` not at pos 2.} \\
\texttt{                        {-} `2` not at pos 3.} \\
\texttt{                    {-} Constraints from `{[}1 2 3 4{]}` (1B, 2W):} \\
\texttt{                        {-} `3` not at pos 4 (from `{[}3 1 2 4{]}`? No, `3` is at pos 1 in guess 3. If `3` was at pos 4, it's W. If `3` was at pos 1, it would be B. But 0B. So `3` not at pos 1).} \\
\texttt{                        {-} We need 1B from `1, 2, 3`.} \\
\texttt{                        {-} `1` cannot be pos 1? (If `1` at pos 1, then in `{[}3 1 2 4{]}`, `1` is at pos 2 {-}> W. `2` at pos 3 {-}> W. `3` at pos 1 {-}> W. `4` out. 3W. Matches).} \\
\texttt{                        {-} `2` cannot be pos 2? (If `2` at pos 2, in `{[}3 1 2 4{]}`, `2` is at pos 3 {-}> W. `1` at pos 2 {-}> B? No, code has `2` at pos 2. Guess has `1` at pos 2. So `1` is W. `3` at pos 1 {-}> W. `4` out. 3W. Matches).} \\
\texttt{                        {-} `3` cannot be pos 3? (If `3` at pos 3, in `{[}3 1 2 4{]}`, `3` is at pos 1 {-}> W. `1` at pos 2 {-}> W. `2` at pos 3 {-}> B? No, code has `3` at pos 3. Guess has `2` at pos 3. So `2` is W. `4` out. 3W. Matches).} \\
\texttt{                    {-} So both cases seem possible so far.} \\
\\ 
\texttt{Let's try to eliminate a digit to see the count change.} \\
\texttt{Guess `{[}1 3 5 6{]}`.} \\
\texttt{{-} If Case A (`3` out, `4` in):} \\
\texttt{    {-} Digits in code: `1, 2, 4, X`.} \\
\texttt{    {-} Guess `{[}1 3 5 6{]}`:} \\
\texttt{        {-} `1` is in.} \\
\texttt{        {-} `3` is out.} \\
\texttt{        {-} `5, 6`? One of them is `X`. So one of `5, 6` is in.} \\
\texttt{        {-} Total correct digits: `1` + `X` = 2.} \\
\texttt{        {-} Feedback should be 0B, 2W (or 1B, 1W).} \\
\texttt{{-} If Case B (`4` out, `3` in):} \\
\texttt{    {-} Digits in code: `1, 2, 3, X`.} \\
\texttt{    {-} Guess `{[}1 3 5 6{]}`:} \\
\texttt{        {-} `1` is in.} \\
\texttt{        {-} `3` is in.} \\
\texttt{        {-} `5, 6`: One is `X`.} \\
\texttt{        {-} Total correct digits: `1` + `3` + `X` = 3.} \\
\texttt{        {-} Feedback should be 0B, 3W (or 1B, 2W).} \\
\\ 
\texttt{This is a great discriminator!} \\
\texttt{{-} If I get 2 pegs, then `3` is out (Case A).} \\
\texttt{{-} If I get 3 pegs, then `3` is in (Case B).} \\
\\ 
\texttt{Let's submit `{[}1 3 5 6{]}`.} \\
\\ 
\texttt{{[}1 3 5 6{]}} \\
            }
        }
    }
    & & \\ \\

    \theutterance \stepcounter{utterance}  
    & & & \multicolumn{4}{p{0.6\linewidth}}{
        \cellcolor[rgb]{0.9,0.9,0.9}{
            \makecell[{{p{\linewidth}}}]{
                \texttt{\tiny{[P1$\langle$GM]}}
                \texttt{{[}GAME{]} Submitted {[}1 3 5 6{]}. Feedback: 0 black peg(s), 2 white peg(s).} \\
            }
        }
    }
     \\ \\

    \theutterance \stepcounter{utterance}  
    & \multicolumn{4}{p{0.6\linewidth}}{
        \cellcolor[rgb]{0.561,0.737,0.561}{
            \makecell[{{p{\linewidth}}}]{
                \texttt{\tiny{[P1$\rangle$GM]}}
                \texttt{Based on the feedback from `{[}1 3 5 6{]}` (0 Black, 2 White):} \\
\texttt{{-} **2 White**: Exactly two digits from my guess are in the code.} \\
\texttt{{-} **Deduction**:} \\
\texttt{    {-} My guess contains `1`, `3`, `5`, `6`.} \\
\texttt{    {-} We previously established that `1` and `2` are definitely in the code.} \\
\texttt{    {-} We are testing if `3` is in the code.} \\
\texttt{    {-} If `3` were in the code, the code would contain `1`, `2`, `3`, and one of `5` or `6` (since we need 4 digits and `4` was the alternative to `3`). That would mean 3 digits from `\{1, 3, 5, 6\}` are in the code (`1`, `3`, and one of `5/6`). This would result in at least 2 or 3 pegs.} \\
\texttt{    {-} However, if `3` is **not** in the code, the code must contain `1`, `2`, `4`, and one of `5` or `6`.} \\
\texttt{        {-} In this scenario (`1, 2, 4, X`), the intersection with `{[}1 3 5 6{]}` is `\{1\}` and `X` (which is either `5` or `6`).} \\
\texttt{        {-} Total correct digits = 2.} \\
\texttt{        {-} This matches the feedback of **2 White pegs**.} \\
\texttt{    {-} Therefore, **`3` is NOT in the code**.} \\
\texttt{    {-} Consequently, **`4` IS in the code**.} \\
\texttt{    {-} The code consists of digits: **`1`, `2`, `4`**, and one of **`5` or `6`**.} \\
\\ 
\texttt{Now we need to determine:} \\
\texttt{1.  Which of `5` or `6` is the fourth digit?} \\
\texttt{2.  What are the positions of `1`, `2`, `4`, and the fifth digit?} \\
\\ 
\texttt{Let's review positional constraints:} \\
\texttt{{-} From `{[}1 2 3 4{]}` (1B, 2W):} \\
\texttt{    {-} `1` is not at pos 1 (because if it were, and `2` is in code, we'd likely see more blacks or the white count would differ? Let's re{-}verify).} \\
\texttt{    {-} Actually, we deduced earlier that `1` and `2` are White in `{[}1 2 3 4{]}` because `3` or `4` was the Black peg.} \\
\texttt{    {-} Wait, if `3` is out, then `4` must be the Black peg in `{[}1 2 3 4{]}`.} \\
\texttt{    {-} So in `{[}1 2 3 4{]}`: `4` is at pos 4 and is the **Black Peg**.} \\
\texttt{    {-} **Conclusion**: The code has **`4` at position 4**.} \\
\\ 
\texttt{Current Code State: `{[}? ? ? 4{]}`} \\
\texttt{Digits remaining to place: `1`, `2`, and one of `\{5, 6\}`.} \\
\texttt{Positions available: 1, 2, 3.} \\
\\ 
\texttt{Let's check the feedback from `{[}3 1 2 4{]}` (0B, 3W):} \\
\texttt{{-} Digits: `3` (out), `1`, `2`, `4` (at pos 4).} \\
\texttt{{-} In the guess `{[}3 1 2 4{]}`:} \\
\texttt{    {-} `4` is at pos 4. Since `4` is at pos 4 in the code, this should be a **Black Peg**.} \\
\texttt{    {-} BUT the feedback was **0 Black**.} \\
\texttt{    {-} **Contradiction!**} \\
\\ 
\texttt{Let's re{-}evaluate the "Black Peg" in `{[}1 2 3 4{]}`.} \\
\texttt{{-} Guess: `1 2 3 4`. Feedback: 1B, 2W.} \\
\texttt{{-} We know `3` is out. So correct digits are `1`, `2`, `4`.} \\
\texttt{{-} If `4` is the Black Peg (pos 4), then in `{[}3 1 2 4{]}`, `4` is at pos 4, so it should be Black. But it's 0B.} \\
\texttt{{-} Therefore, `4` is **NOT** the Black Peg in `{[}1 2 3 4{]}`.} \\
\texttt{{-} So either `1` or `2` is the Black Peg in `{[}1 2 3 4{]}`.} \\
\texttt{{-} If `1` is Black (pos 1):} \\
\texttt{    {-} Code: `1` at pos 1.} \\
\texttt{    {-} Check `{[}3 1 2 4{]}`: `1` is at pos 2. Not Black. `2` at pos 3. `4` at pos 4. `3` out.} \\
\texttt{    {-} If `1` is at pos 1, then in `{[}3 1 2 4{]}`, `1` is W (pos 2). `2` is W (pos 3). `4` is W (pos 4). Total 3W. Matches!} \\
\texttt{    {-} So `1` at pos 1 is a strong candidate.} \\
\texttt{{-} If `2` is Black (pos 2):} \\
\texttt{    {-} Code: `2` at pos 2.} \\
\texttt{    {-} Check `{[}3 1 2 4{]}`: `2` is at pos 3. Not Black. `1` at pos 2. `4` at pos 4. `3` out.} \\
\texttt{    {-} If `2` is at pos 2, then in `{[}3 1 2 4{]}`, `2` is W (pos 3). `1` is W (pos 2). `4` is W (pos 4). Total 3W. Matches!} \\
\texttt{    {-} So `2` at pos 2 is also a candidate.} \\
\\ 
\texttt{So we have two possibilities for the Black Peg in Guess 1:} \\
\texttt{A) Code starts with `1`. (`1` at pos 1).} \\
\texttt{B) Code has `2` at pos 2.} \\
\\ 
\texttt{Let's test these against `{[}1 3 5 6{]}` (0B, 2W).} \\
\texttt{{-} Digits in guess: `1`, `3` (out), `5`, `6`.} \\
\texttt{{-} Correct digits in code: `1`, `2`, `4`, `X` (5 or 6).} \\
\texttt{{-} Intersection: `1` and `X`.} \\
\texttt{{-} Feedback: 2W.} \\
\texttt{{-} This means `1` is White (not at pos 1) AND `X` is White (not at pos 3 or 4).} \\
\texttt{{-} **CRITICAL**: If `1` is White in `{[}1 3 5 6{]}`, then `1` is **NOT** at pos 1.} \\
\texttt{{-} This eliminates Possibility A (`1` at pos 1).} \\
\texttt{{-} Therefore, **`1` is NOT at pos 1**.} \\
\texttt{{-} This implies the Black Peg in `{[}1 2 3 4{]}` was **`2` at pos 2**.} \\
\\ 
\texttt{Current Code State: `{[}? 2 ? 4{]}`} \\
\texttt{Digits remaining: `1`, `X` (5 or 6).} \\
\texttt{Positions available: 1, 3.} \\
\texttt{We know `1` is not at pos 1 (from `{[}1 3 5 6{]}` feedback).} \\
\texttt{So **`1` must be at pos 3**.} \\
\texttt{This leaves **`X` (5 or 6) at pos 1**.} \\
\\ 
\texttt{Current Code State: `{[}X 2 1 4{]}` where `X` is 5 or 6.} \\
\texttt{We need to determine if `X` is 5 or 6.} \\
\texttt{Let's look at `{[}1 3 5 6{]}` again.} \\
\texttt{{-} Guess: `1 3 5 6`.} \\
\texttt{{-} Code: `X 2 1 4`.} \\
\texttt{{-} `1` is at pos 3 in code, pos 1 in guess {-}> White.} \\
\texttt{{-} `X` is at pos 1 in code.} \\
\texttt{    {-} If `X=5`: Guess has `5` at pos 3. Code has `5` at pos 1. {-}> White.} \\
\texttt{    {-} If `X=6`: Guess has `6` at pos 4. Code has `6` at pos 1. {-}> White.} \\
\texttt{{-} `3` is out.} \\
\texttt{{-} Total 2W. This doesn't distinguish 5 vs 6.} \\
\\ 
\texttt{We need a guess that differentiates 5 and 6 or confirms the position of `X`.} \\
\texttt{Let's try to guess `{[}5 2 1 4{]}`.} \\
\texttt{{-} If `X=5`: This is the code. Feedback: 4 Black.} \\
\texttt{{-} If `X=6`: Code is `6 2 1 4`. Guess `5 2 1 4`.} \\
\texttt{    {-} `5` (out). `2` (B). `1` (B). `4` (B).} \\
\texttt{    {-} Feedback: 3 Black.} \\
\\ 
\texttt{Let's submit `{[}5 2 1 4{]}`.} \\
\\ 
\texttt{{[}5 2 1 4{]}} \\
            }
        }
    }
    & & \\ \\

    \theutterance \stepcounter{utterance}  
    & & & \multicolumn{4}{p{0.6\linewidth}}{
        \cellcolor[rgb]{0.9,0.9,0.9}{
            \makecell[{{p{\linewidth}}}]{
                \texttt{\tiny{[P1$\langle$GM]}}
                \texttt{{[}GAME{]} Submitted {[}5 2 1 4{]}. Feedback: 0 black peg(s), 2 white peg(s).} \\
            }
        }
    }
     \\ \\

    \theutterance \stepcounter{utterance}  
    & \multicolumn{4}{p{0.6\linewidth}}{
        \cellcolor[rgb]{0.561,0.737,0.561}{
            \makecell[{{p{\linewidth}}}]{
                \texttt{\tiny{[P1$\rangle$GM]}}
                \texttt{Based on the feedback from `{[}5 2 1 4{]}` (0 Black, 2 White):} \\
\texttt{{-} **0 Black**: None of the digits are in the correct position.} \\
\texttt{    {-} `5` is not at pos 1.} \\
\texttt{    {-} `2` is not at pos 2.} \\
\texttt{    {-} `1` is not at pos 3.} \\
\texttt{    {-} `4` is not at pos 4.} \\
\texttt{{-} **2 White**: Two digits from my guess are in the code but in the wrong positions.} \\
\texttt{{-} **Deduction**:} \\
\texttt{    {-} My previous deduction that the code is `{[}X 2 1 4{]}` relied on the assumption that `2` was the black peg in `{[}1 2 3 4{]}` and `1` was not at pos 1.} \\
\texttt{    {-} However, the feedback `0 Black` for `{[}5 2 1 4{]}` immediately invalidates the positions `2` at pos 2, `1` at pos 3, and `4` at pos 4.} \\
\texttt{    {-} This means my previous positional deductions were incorrect. I need to restart the positional analysis with the confirmed set of digits: **\{1, 2, 4, and one of 5/6\}**.} \\
\texttt{    {-} Let's re{-}evaluate the Black Peg in `{[}1 2 3 4{]}` (1B, 2W).} \\
\texttt{        {-} Digits in code: `1, 2, 4, X`.} \\
\texttt{        {-} If `1` is Black (pos 1): Code starts with `1`.} \\
\texttt{            {-} Check `{[}3 1 2 4{]}` (0B, 3W): `1` is at pos 2 (W). `2` at pos 3 (W). `4` at pos 4 (W). `3` out. Total 3W. Consistent.} \\
\texttt{            {-} Check `{[}1 3 5 6{]}` (0B, 2W): `1` is at pos 1. If code starts with `1`, this should be **Black**. But feedback is 0B.} \\
\texttt{            {-} **Contradiction**: `1` cannot be at pos 1.} \\
\texttt{        {-} If `2` is Black (pos 2): Code has `2` at pos 2.} \\
\texttt{            {-} Check `{[}3 1 2 4{]}` (0B, 3W): `2` is at pos 3 (W). `1` at pos 2 (W). `4` at pos 4 (W). `3` out. Total 3W. Consistent.} \\
\texttt{            {-} Check `{[}1 3 5 6{]}` (0B, 2W): `1` is at pos 1. If code has `2` at pos 2, `1` must be elsewhere. In `{[}1 3 5 6{]}`, `1` is at pos 1. If `1` is not at pos 1 in code, it's White. `X` (5 or 6) is at pos 3 or 4 in guess. If `X` is in code, it's White. Total 2W. Consistent.} \\
\texttt{            {-} So `2` at pos 2 is still a strong candidate.} \\
\texttt{        {-} If `4` is Black (pos 4): Code has `4` at pos 4.} \\
\texttt{            {-} Check `{[}3 1 2 4{]}` (0B, 3W): `4` is at pos 4. Should be Black. But feedback is 0B.} \\
\texttt{            {-} **Contradiction**: `4` cannot be at pos 4.} \\
\texttt{        {-} If `X` (5 or 6) is the Black Peg in `{[}1 2 3 4{]}`:} \\
\texttt{            {-} Then `1, 2, 4` are White.} \\
\texttt{            {-} `1` not at pos 1, `2` not at pos 2, `4` not at pos 4.} \\
\texttt{            {-} Check `{[}3 1 2 4{]}` (0B, 3W): `1` at pos 2 (W), `2` at pos 3 (W), `4` at pos 4 (W). Consistent.} \\
\texttt{            {-} Check `{[}1 3 5 6{]}` (0B, 2W): `1` at pos 1 (W). `X` is in code. If `X` is 5, it's at pos 3 in guess. If `X` is 6, it's at pos 4 in guess.} \\
\texttt{            {-} This scenario is also possible.} \\
\\ 
\texttt{    {-} Let's look at `{[}5 2 1 4{]}` (0B, 2W) again.} \\
\texttt{        {-} Digits: `5, 2, 1, 4`.} \\
\texttt{        {-} We know `2` is not at pos 2 (from 0B).} \\
\texttt{        {-} We know `1` is not at pos 3 (from 0B).} \\
\texttt{        {-} We know `4` is not at pos 4 (from 0B).} \\
\texttt{        {-} We know `5` is not at pos 1 (from 0B).} \\
\texttt{        {-} The 2 White pegs must be from `\{2, 1, 4\}` (since `5` is likely the non{-}code digit or the other of 5/6).} \\
\texttt{        {-} If `5` is the non{-}code digit, then `2, 1, 4` are in the code.} \\
\texttt{            {-} `2` is in code, not at pos 2.} \\
\texttt{            {-} `1` is in code, not at pos 3.} \\
\texttt{            {-} `4` is in code, not at pos 4.} \\
\texttt{        {-} If `5` is in the code, then one of `2, 1, 4` is NOT in the code.} \\
\texttt{            {-} But we established earlier that `1, 2, 4` are almost certainly in the code.} \\
\texttt{            {-} Let's assume the code is a permutation of `1, 2, 4, 5` or `1, 2, 4, 6`.} \\
\texttt{            {-} In `{[}5 2 1 4{]}`, if `5` is in the code, it contributes 1 peg. Then one of `2, 1, 4` contributes 1 peg.} \\
\texttt{            {-} If `5` is NOT in the code, then `2, 1, 4` contribute 2 pegs.} \\
\\ 
\texttt{    {-} Let's try to fix the position of `2`.} \\
\texttt{        {-} In `{[}1 2 3 4{]}`, if `2` was Black, it's at pos 2. But `{[}5 2 1 4{]}` says `2` is not at pos 2.} \\
\texttt{        {-} So `2` is NOT at pos 2.} \\
\texttt{        {-} In `{[}1 2 3 4{]}`, if `2` was White, it's not at pos 2.} \\
\texttt{        {-} So `2` is definitely not at pos 2.} \\
\texttt{        {-} Where can `2` be? Pos 1, 3, or 4.} \\
\texttt{        {-} In `{[}5 2 1 4{]}`, `2` is at pos 2 (White). So `2` is in the code.} \\
\texttt{        {-} In `{[}1 2 3 4{]}`, `2` is at pos 2. If `2` is White there, it's consistent.} \\
\texttt{        {-} If `2` is the Black Peg in `{[}1 2 3 4{]}`, it must be at pos 2. But we just proved `2` is not at pos 2.} \\
\texttt{        {-} **Therefore, `2` is NOT the Black Peg in `{[}1 2 3 4{]}`.**} \\
\texttt{        {-} This eliminates the "2 at pos 2" hypothesis.} \\
\texttt{        {-} So the Black Peg in `{[}1 2 3 4{]}` must be `1` (at pos 1) or `4` (at pos 4) or `X` (5/6).} \\
\texttt{        {-} We already ruled out `1` at pos 1 (due to `{[}1 3 5 6{]}` giving 0B).} \\
\texttt{        {-} We already ruled out `4` at pos 4 (due to `{[}3 1 2 4{]}` giving 0B).} \\
\texttt{        {-} **Therefore, the Black Peg in `{[}1 2 3 4{]}` MUST be `X` (either 5 or 6).**} \\
\texttt{        {-} This implies `X` is at pos 4 in the code.} \\
\texttt{        {-} And `1, 2, 4` are White pegs in `{[}1 2 3 4{]}`.} \\
\texttt{            {-} `1` not at pos 1.} \\
\texttt{            {-} `2` not at pos 2.} \\
\texttt{            {-} `4` not at pos 4.} \\
\\ 
\texttt{    {-} Current State:} \\
\texttt{        {-} Code contains `1, 2, 4` and `X` (5 or 6).} \\
\texttt{        {-} `X` is at **pos 4**.} \\
\texttt{        {-} `1` is not at pos 1.} \\
\texttt{        {-} `2` is not at pos 2.} \\
\texttt{        {-} `4` is not at pos 4.} \\
\texttt{        {-} From `{[}5 2 1 4{]}` (0B, 2W):} \\
\texttt{            {-} `X` is at pos 4 in code. In guess, `4` is at pos 4. `X` is at pos 1 (if 5) or pos 4 (if 6).} \\
\texttt{            {-} If `X=5`: Guess has `5` at pos 1. Code has `5` at pos 4. {-}> White.} \\
\texttt{            {-} If `X=6`: Guess has `6` at pos 4. Code has `6` at pos 4. {-}> **Black**.} \\
\texttt{            {-} But feedback is 0B. So `X` cannot be 6.} \\
\texttt{            {-} **Therefore, `X` is 5.**} \\
\texttt{            {-} And `5` is at pos 4.} \\
\texttt{            {-} So the code ends in `5`.} \\
\texttt{            {-} In `{[}5 2 1 4{]}`, `5` is at pos 1. Code has `5` at pos 4. {-}> White.} \\
\texttt{            {-} We need 1 more White peg from `\{2, 1, 4\}`.} \\
\texttt{            {-} `2` is at pos 2 in guess. Code has `2` not at pos 2. {-}> White.} \\
\texttt{            {-} `1` is at pos 3 in guess. Code has `1` not at pos 3. {-}> White.} \\
\texttt{            {-} `4` is at pos 4 in guess. Code has `4` not at pos 4. {-}> White.} \\
\texttt{            {-} Wait, if `2, 1, 4` are all in the code, and none are in the correct position in `{[}5 2 1 4{]}`, that would be 3 White pegs + 1 White peg for `5` = 4 White pegs.} \\
\texttt{            {-} But we got 2 White pegs.} \\
\texttt{            {-} This implies one of `2, 1, 4` is **NOT** in the code.} \\
\texttt{            {-} But we deduced `1, 2, 4` are in the code because `X` was the Black peg in `{[}1 2 3 4{]}` (contributing 1) and we needed 2 more pegs.} \\
\texttt{            {-} Let's re{-}read `{[}1 2 3 4{]}` (1B, 2W).} \\
\texttt{                {-} If `X` is Black (pos 4), then `1, 2, 3` are the other digits. `3` is out. So `1, 2` are White. That's 1B + 2W. Correct.} \\
\texttt{                {-} So `1` and `2` are definitely in the code.} \\
\texttt{                {-} What about `4`?} \\
\texttt{                {-} If `4` is in the code, then in `{[}1 2 3 4{]}`, we have `1, 2, 4` in code. `X` is also in code? No, code has 4 digits.} \\
\texttt{                {-} If `X` is in code (pos 4), and `1, 2` are in code, that's 3 digits. The 4th digit must be `4`?} \\
\texttt{                {-} If the code is `1, 2, 4, 5`.} \\
\texttt{                {-} Guess `{[}1 2 3 4{]}`: `1` (W), `2` (W), `3` (out), `4` (W? No, `4` is in code).} \\
\texttt{                {-} If code is `1, 2, 4, 5`:} \\
\texttt{                    {-} `1` at pos 1? No, `1` is W.} \\
\texttt{                    {-} `2` at pos 2? No, `2` is W.} \\
\texttt{                    {-} `4` at pos 4? No, `5` is at pos 4. `4` is at pos 3?} \\
\texttt{                    {-} If code is `4 2 1 5`:} \\
\texttt{                        {-} `{[}1 2 3 4{]}`: `1` (W, pos 3), `2` (B, pos 2). Wait, `2` is at pos 2 in guess. If code has `2` at pos 2, it's Black.} \\
\texttt{                        {-} We need 1B. If `2` is B, then `1` and `4` are W.} \\
\texttt{                        {-} Check `{[}5 2 1 4{]}` with `4 2 1 5`:} \\
\texttt{                            {-} `5` (B, pos 4). But feedback is 0B.} \\
\texttt{                            {-} So `5` cannot be at pos 4 if `5` is in the code?} \\
\texttt{                            {-} Wait, if code is `4 2 1 5`, then `5` is at pos 4.} \\
\texttt{                            {-} Guess `{[}5 2 1 4{]}`: `5` at pos 1 (W), `2` at pos 2 (B), `1` at pos 3 (B), `4` at pos 4 (W).} \\
\texttt{                            {-} Result: 2B, 2W.} \\
\texttt{                            {-} Feedback was 0B, 2W.} \\
\texttt{                            {-} So `4 2 1 5` is wrong.} \\
\\ 
\texttt{    {-} Let's reconsider the set of digits.} \\
\texttt{    {-} Maybe `4` is NOT in the code?} \\
\texttt{    {-} If code is `1, 2, 5, 6`?} \\
\texttt{        {-} `{[}1 2 3 4{]}`: `1, 2` in. `3, 4` out. 2 pegs. But we got 3. So `4` must be in or `3` must be in.} \\
\texttt{    {-} Maybe `1` or `2` is out?} \\
\texttt{        {-} If `1` out: `{[}1 2 3 4{]}` has `2, 3, 4`. 3 pegs. `{[}1 3 5 6{]}` has `3, 5, 6`. If `3` in, `5/6` in.} \\
\texttt{        {-} If `1` out, `2` in, `3` in, `4` in, `5` in?} \\
\texttt{        {-} Let's go back to the most solid deduction: **`X` (5 or 6) is the Black Peg in `{[}1 2 3 4{]}` at pos 4.**} \\
\texttt{        {-} This means `5` or `6` is at pos 4.} \\
\texttt{        {-} And `1, 2` are White in `{[}1 2 3 4{]}`.} \\
\texttt{        {-} And `3` is out.} \\
\texttt{        {-} And `4`? If `4` is in the code, we have `1, 2, 4, X`.} \\
\texttt{        {-} If `4` is out, we have `1, 2, 3, X`? No, `3` is out.} \\
\texttt{        {-} So `4` MUST be in the code.} \\
\texttt{        {-} So code is `1, 2, 4, X`.} \\
\texttt{        {-} Why did `{[}5 2 1 4{]}` give 0B, 2W?} \\
\texttt{            {-} `X` is at pos 4.} \\
\texttt{            {-} If `X=5`: Guess `5` at pos 1. Code `5` at pos 4. {-}> W.} \\
\texttt{            {-} If `X=6`: Guess `6` at pos 4. Code `6` at pos 4. {-}> B. (Ruled out).} \\
\texttt{            {-} So `X=5`. Code ends in `5`.} \\
\texttt{            {-} Guess `{[}5 2 1 4{]}`: `5` (W).} \\
\texttt{            {-} We need 1 more W from `2, 1, 4`.} \\
\texttt{            {-} `2` is at pos 2 in guess. Code `2` not at pos 2. {-}> W.} \\
\texttt{            {-} `1` is at pos 3 in guess. Code `1` not at pos 3. {-}> W.} \\
\texttt{            {-} `4` is at pos 4 in guess. Code `4` not at pos 4. {-}> W.} \\
\texttt{            {-} This gives 3W + 1W = 4W.} \\
\texttt{            {-} But we got 2W.} \\
\texttt{            {-} This implies **only one** of `2, 1, 4` is in the code?} \\
\texttt{            {-} But we know `1` and `2` are in the code (from `{[}1 2 3 4{]}` having 3 pegs with `3` out).} \\
\texttt{            {-} Contradiction.} \\
\\ 
\texttt{    {-} Let's rethink: Is `3` in the code?} \\
\texttt{    {-} If `3` is in the code, then `4` is out.} \\
\texttt{    {-} Code: `1, 2, 3, X`.} \\
\texttt{    {-} `{[}1 2 3 4{]}`: `1, 2, 3` in. `4` out. 3 pegs. (1B, 2W).} \\
\texttt{    {-} `{[}3 1 2 4{]}`: `3, 1, 2` in. `4` out. 3 pegs. (0B, 3W).} \\
\texttt{    {-} `{[}1 3 5 6{]}`: `1, 3` in. `5, 6`? One is `X`. So 3 digits in.} \\
\texttt{        {-} Feedback was 2W.} \\
\texttt{        {-} If 3 digits in, we should get 2 or 3 pegs. 2W is possible if 1 is B? No, 0B. So 2W means 2 digits in?} \\
\texttt{        {-} If `1, 3, X` are in, that's 3 digits. 0B means all wrong pos. Should be 3W.} \\
\texttt{        {-} But we got 2W.} \\
\texttt{        {-} This implies only **2** digits from `1, 3, 5, 6` are in the code.} \\
\texttt{        {-} We know `1` is in. `3` is in. That's 2.} \\
\texttt{        {-} So `X` (5 or 6) is **NOT** in the code.} \\
\texttt{        {-} So the code is `1, 2, 3, 4`? No, `4` is out in this hypothesis.} \\
\texttt{        {-} So code is `1, 2, 3, Y` where `Y` is not 5 or 6?} \\
\texttt{        {-} Digits 1{-}6. If `1, 2, 3` are in, and `5, 6` are out, then `4` must be in.} \\
\texttt{        {-} But we assumed `4` is out.} \\
\texttt{        {-} Contradiction.} \\
\\ 
\texttt{    {-} Let's try: **`1` is NOT in the code.**} \\
\texttt{    {-} Code: `2, 3, 4, X`.} \\
\texttt{    {-} `{[}1 2 3 4{]}`: `2, 3, 4` in. 3 pegs. (1B, 2W).} \\
\texttt{    {-} `{[}3 1 2 4{]}`: `3, 2, 4` in. 3 pegs. (0B, 3W).} \\
\texttt{    {-} `{[}1 3 5 6{]}`: `3` in. `5, 6`? One is `X`. So 2 digits in.} \\
\texttt{        {-} Feedback 2W. Matches!} \\
\texttt{    {-} So `1` is out. `X` is in (5 or 6).} \\
\texttt{    {-} Code: `2, 3, 4, X`.} \\
\texttt{    {-} `{[}1 2 3 4{]}` (1B, 2W):} \\
\texttt{        {-} `2` at pos 2? If B, then `3, 4` W.} \\
\texttt{        {-} `3` at pos 3? If B, then `2, 4` W.} \\
\texttt{        {-} `4` at pos 4? If B, then `2, 3` W.} \\
\texttt{    {-} `{[}3 1 2 4{]}` (0B, 3W):} \\
\texttt{        {-} `3` at pos 1 (W).} \\
\texttt{        {-} `2` at pos 3 (W).} \\
\texttt{        {-} `4` at pos 4 (W).} \\
\texttt{        {-} So `3` not at 1, `2` not at 3, `4` not at 4.} \\
\texttt{    {-} `{[}5 2 1 4{]}` (0B, 2W):} \\
\texttt{        {-} `2` at pos 2 (W). So `2` not at 2.} \\
\texttt{        {-} `4` at pos 4 (W). So `4` not at 4.} \\
\texttt{        {-} `5` at pos 1. If `X=5`, then `5` is in code. `5` at pos 1 (W) {-}> `5` not at 1.} \\
\texttt{        {-} `1` is out.} \\
\texttt{        {-} We need 2W. `2` and `4` are W. That's 2.} \\
\texttt{        {-} So `5` must be out? Or `X` is not 5?} \\
\texttt{        {-} If `X=6`: `5` is out. `2, 4` are W. `6` is not in guess. Total 2W. Matches!} \\
\texttt{        {-} So `X=6`.} \\
\texttt{        {-} Code: `2, 3, 4, 6`.} \\
\texttt{        {-} Positions:} \\
\texttt{            {-} `2` not at 2, 3.} \\
\texttt{            {-} `3` not at 1.} \\
\texttt{            {-} `4` not at 4.} \\
\texttt{            {-} `6` not at 4 (from `{[}1 2 3 4{]}} \\
            }
        }
    }
    & & \\ \\

    \theutterance \stepcounter{utterance}  
    & & & \multicolumn{4}{p{0.6\linewidth}}{
        \cellcolor[rgb]{0.9,0.9,0.9}{
            \makecell[{{p{\linewidth}}}]{
                \texttt{\tiny{[P1$\langle$GM]}}
                \texttt{{[}GAME{]} You attempted an invalid move. Reason: You have already guessed {[}1, 2, 3, 4{]}. Please try a different guess. Please resubmit a valid move and remember to follow the game rules to avoid penalties.} \\
            }
        }
    }
     \\ \\

    \theutterance \stepcounter{utterance}  
    & \multicolumn{4}{p{0.6\linewidth}}{
        \cellcolor[rgb]{0.561,0.737,0.561}{
            \makecell[{{p{\linewidth}}}]{
                \texttt{\tiny{[P1$\rangle$GM]}}
                \texttt{{[}2 4 1 5{]}} \\
            }
        }
    }
    & & \\ \\

    \theutterance \stepcounter{utterance}  
    & & & \multicolumn{4}{p{0.6\linewidth}}{
        \cellcolor[rgb]{0.9,0.9,0.9}{
            \makecell[{{p{\linewidth}}}]{
                \texttt{\tiny{[P1$\langle$GM]}}
                \texttt{{[}GAME{]} Submitted {[}2 4 1 5{]}. Feedback: 1 black peg(s), 1 white peg(s).} \\
            }
        }
    }
     \\ \\

    \theutterance \stepcounter{utterance}  
    & \multicolumn{4}{p{0.6\linewidth}}{
        \cellcolor[rgb]{0.561,0.737,0.561}{
            \makecell[{{p{\linewidth}}}]{
                \texttt{\tiny{[P1$\rangle$GM]}}
                \texttt{Based on the feedback from `{[}2 4 1 5{]}` (1 Black, 1 White):} \\
\texttt{{-} **1 Black**: One digit is correct and in the correct position.} \\
\texttt{{-} **1 White**: One digit is correct but in the wrong position.} \\
\texttt{{-} **Total correct digits**: 2.} \\
\\ 
\texttt{Let's re{-}evaluate the set of digits in the code.} \\
\texttt{Previous guesses:} \\
\texttt{1. `{[}1 2 3 4{]}` {-}> 1B, 2W (3 correct digits)} \\
\texttt{2. `{[}1 2 5 6{]}` {-}> 0B, 2W (2 correct digits)} \\
\texttt{3. `{[}3 1 2 4{]}` {-}> 0B, 3W (3 correct digits)} \\
\texttt{4. `{[}1 3 5 6{]}` {-}> 0B, 2W (2 correct digits)} \\
\texttt{5. `{[}2 4 1 5{]}` {-}> 1B, 1W (2 correct digits)} \\
\\ 
\texttt{**Deduction on Digits:**} \\
\texttt{{-} From `{[}1 2 5 6{]}` (2 correct) and `{[}1 3 5 6{]}` (2 correct):} \\
\texttt{    {-} Common digits: `1`, `5`, `6`.} \\
\texttt{    {-} If `1` is in the code, then `5` and `6` must be out (otherwise we'd have 3 correct digits in `{[}1 3 5 6{]}`).} \\
\texttt{    {-} If `1` is out, then `5` and `6` must be in (to get 2 pegs in both).} \\
\texttt{    {-} Let's test **Case A: `1` is in, `5` and `6` are out.**} \\
\texttt{        {-} Code contains `1`.} \\
\texttt{        {-} From `{[}1 2 5 6{]}` (2 correct): `1` and one of `2` are in. So `2` is in.} \\
\texttt{        {-} From `{[}1 3 5 6{]}` (2 correct): `1` and one of `3` are in. So `3` is in.} \\
\texttt{        {-} So far: `1, 2, 3` are in.} \\
\texttt{        {-} From `{[}1 2 3 4{]}` (3 correct): `1, 2, 3` are in. `4` is out.} \\
\texttt{        {-} So the code is a permutation of **`1, 2, 3` and one other digit**.} \\
\texttt{        {-} The other digit must be from `\{4, 5, 6\}`. Since `4` is out (from `{[}1 2 3 4{]}` having 3 correct with `1,2,3`), and `5, 6` are out (from `{[}1 3 5 6{]}`), this creates a contradiction. We need 4 digits.} \\
\texttt{        {-} Wait, if `1, 2, 3` are in, that's 3 digits. The 4th digit must be `4`, `5`, or `6`.} \\
\texttt{        {-} If `4` is in: `{[}1 2 3 4{]}` has 4 correct digits. But we got 3 pegs. So `4` is out.} \\
\texttt{        {-} If `5` is in: `{[}1 3 5 6{]}` has `1, 3, 5` in. That's 3 correct digits. But we got 2 pegs. So `5` is out.} \\
\texttt{        {-} If `6` is in: `{[}1 3 5 6{]}` has `1, 3, 6` in. That's 3 correct digits. But we got 2 pegs. So `6` is out.} \\
\texttt{        {-} **Conclusion**: Case A (`1` is in) leads to a contradiction where we can't find a 4th digit.} \\
\\ 
\texttt{{-} Let's test **Case B: `1` is out, `5` and `6` are in.**} \\
\texttt{    {-} Code contains `5`, `6`.} \\
\texttt{    {-} From `{[}1 2 5 6{]}` (2 correct): `5, 6` are in. `1, 2` are out. So `2` is out.} \\
\texttt{    {-} From `{[}1 3 5 6{]}` (2 correct): `5, 6` are in. `1, 3` are out. So `3` is out.} \\
\texttt{    {-} So far: `5, 6` are in. `1, 2, 3` are out.} \\
\texttt{    {-} Remaining digits: `4`. We need 2 more digits to make 4. But we only have `4` left? No, digits are 1{-}6.} \\
\texttt{    {-} If `1, 2, 3` are out, the code must be `4, 5, 6` and... wait, we need 4 digits.} \\
\texttt{    {-} This implies my deduction about "out" is too strict. Let's re{-}evaluate the counts.} \\
\texttt{    {-} Maybe `1` is in, but `5` or `6` is also in?} \\
\texttt{    {-} Let's look at `{[}2 4 1 5{]}` (1B, 1W).} \\
\texttt{        {-} Digits: `2, 4, 1, 5`.} \\
\texttt{        {-} If `1` is in:} \\
\texttt{            {-} If `1` is the Black peg (pos 3): Code has `1` at pos 3.} \\
\texttt{            {-} Then we need 1 White peg from `2, 4, 5`.} \\
\texttt{            {-} If `1` is in, and we established earlier that `1, 2, 3` being in causes issues with `{[}1 3 5 6{]}`, maybe `2` or `3` is out?} \\
\texttt{            {-} Let's try a different angle.} \\
\texttt{            {-} Look at `{[}1 2 3 4{]}` (1B, 2W) and `{[}2 4 1 5{]}` (1B, 1W).} \\
\texttt{            {-} Common digits: `1, 2, 4`.} \\
\texttt{            {-} In `{[}1 2 3 4{]}`, `1` is at pos 1, `2` at pos 2, `4` at pos 4.} \\
\texttt{            {-} In `{[}2 4 1 5{]}`, `2` is at pos 1, `4` at pos 2, `1` at pos 3.} \\
\texttt{            {-} If `1` is the Black peg in `{[}1 2 3 4{]}` (pos 1):} \\
\texttt{                {-} In `{[}2 4 1 5{]}`, `1` is at pos 3. It should be White.} \\
\texttt{                {-} We have 1B, 1W. So `1` (W) accounts for the White peg.} \\
\texttt{                {-} Then `2, 4, 5` must be out.} \\
\texttt{                {-} If `2, 4, 5` are out, code has `1` and `3, 6`?} \\
\texttt{                {-} Check `{[}1 2 3 4{]}`: `1` (B), `2` (out), `3` (W?), `4` (out).} \\
\texttt{                {-} If `3` is in, it's at pos 3 in guess. If code has `3` elsewhere, it's W.} \\
\texttt{                {-} We need 2W. `1` is B. `3` is W. Where is the other W?} \\
\texttt{                {-} Maybe `4` is in? If `4` is in, `{[}2 4 1 5{]}` has `4` (W or B).} \\
\texttt{                {-} This path is getting complicated.} \\
\\ 
\texttt{Let's try a guess that isolates the Black peg position more clearly.} \\
\texttt{We know the code has 4 distinct digits from 1{-}6.} \\
\texttt{Let's guess `{[}4 5 6 1{]}`.} \\
\texttt{{-} If `1` is in the code:} \\
\texttt{    {-} If `1` is the Black peg in `{[}1 2 3 4{]}` (pos 1), then in `{[}4 5 6 1{]}`, `1` is at pos 4 (White).} \\
\texttt{    {-} If `1` is the Black peg in `{[}2 4 1 5{]}` (pos 3), then in `{[}4 5 6 1{]}`, `1` is at pos 4 (White).} \\
\texttt{{-} Let's try to confirm if `4` is in the code.} \\
\texttt{{-} Guess `{[}4 1 2 3{]}`.} \\
\texttt{    {-} If `4` is in the code:} \\
\texttt{        {-} In `{[}1 2 3 4{]}`, `4` is at pos 4.} \\
\texttt{        {-} In `{[}2 4 1 5{]}`, `4` is at pos 2.} \\
\texttt{        {-} In `{[}4 1 2 3{]}`, `4` is at pos 1.} \\
\texttt{    {-} If `4` is the Black peg in `{[}1 2 3 4{]}` (pos 4):} \\
\texttt{        {-} In `{[}2 4 1 5{]}`, `4` is at pos 2 (White).} \\
\texttt{        {-} In `{[}4 1 2 3{]}`, `4` is at pos 1 (White).} \\
\texttt{    {-} If `4` is NOT in the code:} \\
\texttt{        {-} In `{[}1 2 3 4{]}`, `4` contributes 0.} \\
\texttt{        {-} In `{[}2 4 1 5{]}`, `4` contributes 0.} \\
\texttt{        {-} Then `{[}1 2 3 4{]}` has `1, 2, 3` in. `{[}2 4 1 5{]}` has `1, 2` in? No, `4` out. `1, 2, 5` in?} \\
\texttt{        {-} If `4` is out, code is subset of `1, 2, 3, 5, 6`.} \\
\texttt{        {-} `{[}1 2 3 4{]}` (3 pegs) {-}> `1, 2, 3` in.} \\
\texttt{        {-} `{[}2 4 1 5{]}` (2 pegs) {-}> `1, 2` in (since `4` out). `5` out?} \\
\texttt{        {-} If `1, 2, 3` in, and `5` out, then `6` must be in.} \\
\texttt{        {-} Code: `1, 2, 3, 6`.} \\
\texttt{        {-} Check `{[}1 2 3 4{]}`: `1, 2, 3` in. `4` out. 3 pegs. (1B, 2W).} \\
\texttt{        {-} Check `{[}2 4 1 5{]}`: `2, 1` in. `4, 5` out. 2 pegs. (1B, 1W).} \\
\texttt{        {-} This fits perfectly!} \\
\texttt{        {-} So the code is a permutation of **`1, 2, 3, 6`**.} \\
\texttt{        {-} `4` and `5` are out.} \\
\\ 
\texttt{Now let's find the positions for `1, 2, 3, 6`.} \\
\texttt{{-} From `{[}1 2 3 4{]}` (1B, 2W):} \\
\texttt{    {-} `1` at pos 1, `2` at pos 2, `3` at pos 3.} \\
\texttt{    {-} One is Black, two are White.} \\
\texttt{{-} From `{[}2 4 1 5{]}` (1B, 1W):} \\
\texttt{    {-} `2` at pos 1, `1` at pos 3. (`4, 5` out).} \\
\texttt{    {-} One is Black, one is White.} \\
\texttt{{-} From `{[}3 1 2 4{]}` (0B, 3W):} \\
\texttt{    {-} `3` at pos 1, `1` at pos 2, `2` at pos 3.} \\
\texttt{    {-} All three (`1, 2, 3`) are White.} \\
\texttt{    {-} So:} \\
\texttt{        {-} `3` is not at pos 1.} \\
\texttt{        {-} `1` is not at pos 2.} \\
\texttt{        {-} `2` is not at pos 3.} \\
\texttt{{-} From `{[}1 3 5 6{]}` (0B, 2W):} \\
\texttt{    {-} `1` at pos 1, `3` at pos 2, `6` at pos 3 or 4.} \\
\texttt{    {-} `1` is White (not at pos 1).} \\
\texttt{    {-} `3` is White (not at pos 2).} \\
\texttt{    {-} `6` is White (not at pos 3 or 4).} \\
\texttt{    {-} So `6` is not at pos 3 or 4. `6` must be at pos 1 or 2.} \\
\\ 
\texttt{Let's combine constraints:} \\
\texttt{{-} `1` not at pos 1 (from `{[}1 3 5 6{]}`).} \\
\texttt{{-} `1` not at pos 2 (from `{[}3 1 2 4{]}`).} \\
\texttt{{-} So **`1` is at pos 3 or 4**.} \\
\texttt{{-} `2` not at pos 2 (from `{[}1 2 3 4{]}`? No, we don't know if `2` is B there yet).} \\
\texttt{{-} `2` not at pos 3 (from `{[}3 1 2 4{]}`).} \\
\texttt{{-} `2` not at pos 1? In `{[}2 4 1 5{]}`, `2` is at pos 1. If `2` was Black there, then `1` is White.} \\
\texttt{    {-} If `2` is Black in `{[}2 4 1 5{]}` (pos 1):} \\
\texttt{        {-} Then `1` is White in `{[}2 4 1 5{]}` (pos 3).} \\
\texttt{        {-} In `{[}1 2 3 4{]}`, `2` is at pos 1. If `2` is Black, then `1` and `3` are White.} \\
\texttt{        {-} In `{[}3 1 2 4{]}`, `2` is at pos 3. If `2` is at pos 1, it's White. `1` at pos 2 (White). `3` at pos 1 (White). Total 3W. Matches.} \\
\texttt{        {-} So `2` at pos 1 is a strong candidate.} \\
\texttt{{-} `3` not at pos 1 (from `{[}3 1 2 4{]}`).} \\
\texttt{{-} `3` not at pos 2 (from `{[}1 3 5 6{]}`).} \\
\texttt{{-} `3` not at pos 3? In `{[}1 2 3 4{]}`, `3` is at pos 3. If `3` is Black, then `1, 2` are White.} \\
\texttt{    {-} If `3` is Black (pos 3):} \\
\texttt{        {-} In `{[}3 1 2 4{]}`, `3` is at pos 1 (White). `1` at pos 2 (White). `2` at pos 3 (White). Total 3W. Matches.} \\
\texttt{        {-} In `{[}1 2 3 4{]}`, `3` is B. `1, 2` are W. Matches.} \\
\texttt{        {-} In `{[}2 4 1 5{]}`, `3` is not in guess. `1` at pos 3 (W). `2` at pos 1 (W). Total 2W. But we got 1B, 1W.} \\
\texttt{        {-} So `3` cannot be the Black peg in `{[}1 2 3 4{]}`.} \\
\texttt{{-} So the Black peg in `{[}1 2 3 4{]}` is either `1` or `2`.} \\
\texttt{    {-} If `1` is Black (pos 1):} \\
\texttt{        {-} Contradicts `1` not at pos 1 (from `{[}1 3 5 6{]}`).} \\
\texttt{    {-} So **`2` must be the Black peg in `{[}1 2 3 4{]}`**.} \\
\texttt{    {-} Therefore, **`2` is at pos 2**.} \\
\\ 
\texttt{Current State: `{[}? 2 ? ?{]}`} \\
\texttt{Digits: `1, 3, 6`.} \\
\texttt{Constraints:} \\
\texttt{{-} `1` not at pos 1, not at pos 2. So `1` is at pos 3 or 4.} \\
\texttt{{-} `3` not at pos 1, not at pos 2. So `3` is at pos 3 or 4.} \\
\texttt{{-} `6` not at pos 3, not at pos 4 (from `{[}1 3 5 6{]}`). So `6` is at pos 1.} \\
\texttt{{-} Therefore, **`6` is at pos 1**.} \\
\\ 
\texttt{Current State: `{[}6 2 ? ?{]}`} \\
\texttt{Remaining positions: 3, 4.} \\
\texttt{Remaining digits: `1, 3`.} \\
\texttt{{-} `1` is at pos 3 or 4.} \\
\texttt{{-} `3` is at pos 3 or 4.} \\
\texttt{{-} Check `{[}2 4 1 5{]}` (1B, 1W):} \\
\texttt{    {-} `2` is at pos 1. Code has `2` at pos 2. So `2` is White.} \\
\texttt{    {-} We need 1 Black. So either `1` or `5` (out) or `4` (out) must be Black? No, `4, 5` are out.} \\
\texttt{    {-} So `1` must be the Black peg in `{[}2 4 1 5{]}`.} \\
\texttt{    {-} In `{[}2 4 1 5{]}`, `1` is at pos 3.} \\
\texttt{    {-} So **`1` is at pos 3**.} \\
\texttt{{-} Therefore, **`3` is at pos 4**.} \\
\\ 
\texttt{Proposed Code: **`6 2 1 3`**.} \\
\\ 
\texttt{Let's verify:} \\
\texttt{{-} `{[}1 2 3 4{]}`: `1`(W), `2`(B), `3`(W), `4`(out). {-}> 1B, 2W. (Matches)} \\
\texttt{{-} `{[}1 2 5 6{]}`: `1`(W), `2`(W), `5`(out), `6`(W). {-}> 0B, 3W.} \\
\texttt{    {-} Wait, my previous feedback for `{[}1 2 5 6{]}` was **0B, 2W**.} \\
\texttt{    {-} My proposed code `6 2 1 3` gives 3W for `{[}1 2 5 6{]}`.} \\
\texttt{    {-} Contradiction.} \\
\\ 
\texttt{Re{-}evaluate `{[}1 2 5 6{]}` feedback: 0B, 2W.} \\
\texttt{{-} Digits in guess: `1, 2, 5, 6`.} \\
\texttt{{-} Digits in code: `1, 2, 3, 6` (based on `4, 5` out).} \\
\texttt{{-} Intersection: `1, 2, 6`. (3 digits).} \\
\texttt{{-} If 3 digits are in the code, we should get 3 pegs (since 0B).} \\
\texttt{{-} But we got 2W.} \\
\texttt{{-} This means one of `1, 2, 6` is **NOT** in the code.} \\
\texttt{{-} Which one?} \\
\texttt{    {-} If `1` is out: Code `2, 3, 6, 4`? No, `4` out. `2, 3, 6, ?`.} \\
\texttt{    {-} If `2` is out: Code `1, 3, 6, 4`? No, `4` out.} \\
\texttt{    {-} If `6` is out: Code `1, 2, 3, 4`? No, `4` out.} \\
\texttt{    {-} Maybe `4` IS in the code?} \\
\texttt{    {-} Let's re{-}examine `{[}1 2 3 4{]}` (1B, 2W).} \\
\texttt{        {-} If `4` is in, and `1, 2, 3` are in, that's 4 digits.} \\
\texttt{        {-} Then `{[}1 2 3 4{]}` has 4 correct digits. But we got 3 pegs.} \\
\texttt{        {-} So one of `1, 2, 3, 4` is out.} \\
\texttt{    {-} If `4` is in, then one of `1, 2, 3` is out.} \\
\texttt{    {-} If `1` is out: Code `2, 3, 4, 6`.} \\
\texttt{        {-} `{[}1 2 3 4{]}`: `2, 3, 4` in. 3 pegs. (1B, 2W).} \\
\texttt{        {-} `{[}1 2 5 6{]}`: `2, 6` in. `1, 5` out. 2 pegs. (0B, 2W). **Matches!**} \\
\texttt{        {-} `{[}3 1 2 4{]}`: `3, 2, 4` in. 3 pegs. (0B, 3W). **Matches!**} \\
\texttt{        {-} `{[}1 3 5 6{]}`: `3, 6` in. `1, 5` out. 2 pegs. (0B, 2W). **Matches!**} \\
\texttt{        {-} `{[}2 4 1 5{]}`: `2, 4` in. `1, 5` out. 2 pegs. (1B, 1W).} \\
\texttt{            {-} We need 1B, 1W.} \\
\texttt{            {-} Digits `2, 4` are in.} \\
\texttt{            {-} If `2` is Black (pos 1 in guess? No, `2` is at pos 1 in guess).} \\
\texttt{            {-} If `2` is at pos 1 in code:} \\
\texttt{                {-} `{[}2 4 1 5{]}`: `2` (B). `4` (W or B). `1` (out). `5` (out).} \\
\texttt{                {-} If `4` is W, we get 1B, 1W.} \\
\texttt{                {-} Check `{[}1 2 3 4{]}`: `2` at pos 2. If `2` is at pos 1, it's W. `3` at pos 3. `4` at pos 4.} \\
\texttt{                {-} If `2` is W, `3` is W, `4` is B (pos 4).} \\
\texttt{                {-} Then `{[}1 2 3 4{]}` is 1B (`4`), 2W (`2, 3`). Matches.} \\
\texttt{                {-} Check `{[}3 1 2 4{]}`: `3` at pos 1 (W). `1` out. `2` at pos 3 (W). `4` at pos 4 (B).} \\
\texttt{                {-} Result: 1B, 2W. But we got 0B, 3W.} \\
\texttt{                {-} So `4` cannot be Black in `{[}3 1 2 4{]}`.} \\
\texttt{                {-} So `} \\
            }
        }
    }
    & & \\ \\

    \theutterance \stepcounter{utterance}  
    & & & \multicolumn{4}{p{0.6\linewidth}}{
        \cellcolor[rgb]{0.9,0.9,0.9}{
            \makecell[{{p{\linewidth}}}]{
                \texttt{\tiny{[P1$\langle$GM]}}
                \texttt{{[}GAME{]} You attempted an invalid move. Reason: You have already guessed {[}3, 1, 2, 4{]}. Please try a different guess. Please resubmit a valid move and remember to follow the game rules to avoid penalties.} \\
            }
        }
    }
     \\ \\

    \theutterance \stepcounter{utterance}  
    & \multicolumn{4}{p{0.6\linewidth}}{
        \cellcolor[rgb]{0.561,0.737,0.561}{
            \makecell[{{p{\linewidth}}}]{
                \texttt{\tiny{[P1$\rangle$GM]}}
                \texttt{{[}6 2 1 3{]}} \\
            }
        }
    }
    & & \\ \\

    \theutterance \stepcounter{utterance}  
    & & & \multicolumn{4}{p{0.6\linewidth}}{
        \cellcolor[rgb]{0.9,0.9,0.9}{
            \makecell[{{p{\linewidth}}}]{
                \texttt{\tiny{[P1$\langle$GM]}}
                \texttt{{[}GAME{]} Submitted {[}6 2 1 3{]}. Feedback: 1 black peg(s), 2 white peg(s).} \\
            }
        }
    }
     \\ \\

    \theutterance \stepcounter{utterance}  
    & \multicolumn{4}{p{0.6\linewidth}}{
        \cellcolor[rgb]{0.561,0.737,0.561}{
            \makecell[{{p{\linewidth}}}]{
                \texttt{\tiny{[P1$\rangle$GM]}}
                \texttt{Based on the feedback from `{[}6 2 1 3{]}` (1 Black, 2 White):} \\
\texttt{{-} **1 Black**: One digit is correct and in the correct position.} \\
\texttt{{-} **2 White**: Two digits are correct but in the wrong positions.} \\
\texttt{{-} **Total correct digits**: 3.} \\
\\ 
\texttt{Let's re{-}evaluate the set of digits.} \\
\texttt{Previous guesses:} \\
\texttt{1. `{[}1 2 3 4{]}` {-}> 1B, 2W (3 correct)} \\
\texttt{2. `{[}1 2 5 6{]}` {-}> 0B, 2W (2 correct)} \\
\texttt{3. `{[}3 1 2 4{]}` {-}> 0B, 3W (3 correct)} \\
\texttt{4. `{[}1 3 5 6{]}` {-}> 0B, 2W (2 correct)} \\
\texttt{5. `{[}2 4 1 5{]}` {-}> 1B, 1W (2 correct)} \\
\texttt{6. `{[}6 2 1 3{]}` {-}> 1B, 2W (3 correct)} \\
\\ 
\texttt{**Deduction on Digits:**} \\
\texttt{{-} From `{[}1 2 5 6{]}` (2 correct) and `{[}1 3 5 6{]}` (2 correct):} \\
\texttt{    {-} Common digits: `1`, `5`, `6`.} \\
\texttt{    {-} If `1` is in the code, then `5` and `6` must be out (otherwise `{[}1 3 5 6{]}` would have 3 correct digits).} \\
\texttt{    {-} If `1` is out, then `5` and `6` must be in (to get 2 pegs in both).} \\
\texttt{    {-} Let's test **Case A: `1` is in, `5` and `6` are out.**} \\
\texttt{        {-} Code contains `1`.} \\
\texttt{        {-} From `{[}1 2 5 6{]}` (2 correct): `1` and one of `2` are in. So `2` is in.} \\
\texttt{        {-} From `{[}1 3 5 6{]}` (2 correct): `1` and one of `3` are in. So `3` is in.} \\
\texttt{        {-} So far: `1, 2, 3` are in.} \\
\texttt{        {-} From `{[}1 2 3 4{]}` (3 correct): `1, 2, 3` are in. `4` is out.} \\
\texttt{        {-} So the code is a permutation of **`1, 2, 3` and one other digit**.} \\
\texttt{        {-} The other digit must be from `\{4, 5, 6\}`. Since `4` is out (from `{[}1 2 3 4{]}` having 3 correct with `1,2,3`), and `5, 6` are out (from `{[}1 3 5 6{]}`), this creates a contradiction. We need 4 digits.} \\
\texttt{        {-} Wait, if `1, 2, 3` are in, that's 3 digits. The 4th digit must be `4`, `5`, or `6`.} \\
\texttt{        {-} If `4` is in: `{[}1 2 3 4{]}` has 4 correct digits. But we got 3 pegs. So `4` is out.} \\
\texttt{        {-} If `5` is in: `{[}1 3 5 6{]}` has `1, 3, 5` in. That's 3 correct digits. But we got 2 pegs. So `5` is out.} \\
\texttt{        {-} If `6` is in: `{[}1 3 5 6{]}` has `1, 3, 6` in. That's 3 correct digits. But we got 2 pegs. So `6` is out.} \\
\texttt{        {-} **Conclusion**: Case A (`1` is in) leads to a contradiction where we can't find a 4th digit.} \\
\\ 
\texttt{{-} Let's test **Case B: `1` is out, `5` and `6` are in.**} \\
\texttt{    {-} Code contains `5`, `6`.} \\
\texttt{    {-} From `{[}1 2 5 6{]}` (2 correct): `5, 6` are in. `1, 2` are out. So `2` is out.} \\
\texttt{    {-} From `{[}1 3 5 6{]}` (2 correct): `5, 6` are in. `1, 3` are out. So `3` is out.} \\
\texttt{    {-} So far: `5, 6` are in. `1, 2, 3` are out.} \\
\texttt{    {-} Remaining digits: `4`. We need 2 more digits to make 4. But we only have `4` left? No, digits are 1{-}6.} \\
\texttt{    {-} If `1, 2, 3` are out, the code must be `4, 5, 6` and... wait, we need 4 digits.} \\
\texttt{    {-} This implies my deduction about "out" is too strict. Let's re{-}evaluate the counts.} \\
\texttt{    {-} Maybe `1` is in, but `5` or `6` is also in?} \\
\texttt{    {-} Let's look at `{[}2 4 1 5{]}` (1B, 1W).} \\
\texttt{        {-} Digits: `2, 4, 1, 5`.} \\
\texttt{        {-} If `1` is in:} \\
\texttt{            {-} If `1` is the Black peg (pos 3): Code has `1` at pos 3.} \\
\texttt{            {-} Then we need 1 White peg from `2, 4, 5`.} \\
\texttt{            {-} If `1` is in, and we established earlier that `1, 2, 3` being in causes issues with `{[}1 3 5 6{]}`, maybe `2` or `3` is out?} \\
\texttt{            {-} Let's try a different angle.} \\
\texttt{            {-} Look at `{[}1 2 3 4{]}` (1B, 2W) and `{[}2 4 1 5{]}` (1B, 1W).} \\
\texttt{            {-} Common digits: `1, 2, 4`.} \\
\texttt{            {-} In `{[}1 2 3 4{]}`, `1` is at pos 1, `2` at pos 2, `4` at pos 4.} \\
\texttt{            {-} In `{[}2 4 1 5{]}`, `2` is at pos 1, `4` at pos 2, `1` at pos 3.} \\
\texttt{            {-} If `1` is the Black peg in `{[}1 2 3 4{]}` (pos 1):} \\
\texttt{                {-} In `{[}2 4 1 5{]}`, `1` is at pos 3. It should be White.} \\
\texttt{                {-} We have 1B, 1W. So `1` (W) accounts for the White peg.} \\
\texttt{                {-} Then `2, 4, 5` must be out.} \\
\texttt{                {-} If `2, 4, 5` are out, code has `1` and `3, 6`?} \\
\texttt{                {-} Check `{[}1 2 3 4{]}`: `1` (B), `2` (out), `3` (W?), `4` (out).} \\
\texttt{                {-} If `3` is in, it's at pos 3 in guess. If code has `3` elsewhere, it's W.} \\
\texttt{                {-} We need 2W. `1` is B. `3` is W. Where is the other W?} \\
\texttt{                {-} Maybe `4` is in? If `4` is in, `{[}2 4 1 5{]}` has `4` (W or B).} \\
\texttt{                {-} This path is getting complicated.} \\
\\ 
\texttt{Let's try a guess that isolates the Black peg position more clearly.} \\
\texttt{We know the code has 4 distinct digits from 1{-}6.} \\
\texttt{Let's guess `{[}4 5 6 1{]}`.} \\
\texttt{{-} If `1` is in the code:} \\
\texttt{    {-} If `1` is the Black peg in `{[}1 2 3 4{]}` (pos 1), then in `{[}4 5 6 1{]}`, `1` is at pos 4 (White).} \\
\texttt{    {-} If `1` is the Black peg in `{[}2 4 1 5{]}` (pos 3), then in `{[}4 5 6 1{]}`, `1` is at pos 4 (White).} \\
\texttt{{-} Let's try to confirm if `4` is in the code.} \\
\texttt{{-} Guess `{[}4 1 2 3{]}`.} \\
\texttt{    {-} If `4` is in the code:} \\
\texttt{        {-} In `{[}1 2 3 4{]}`, `4` is at pos 4.} \\
\texttt{        {-} In `{[}2 4 1 5{]}`, `4` is at pos 2.} \\
\texttt{        {-} In `{[}4 1 2 3{]}`, `4` is at pos 1.} \\
\texttt{    {-} If `4` is the Black peg in `{[}1 2 3 4{]}` (pos 4):} \\
\texttt{        {-} In `{[}2 4 1 5{]}`, `4` is at pos 2 (White).} \\
\texttt{        {-} In `{[}4 1 2 3{]}`, `4` is at pos 1 (White).} \\
\texttt{    {-} If `4` is NOT in the code:} \\
\texttt{        {-} In `{[}1 2 3 4{]}`, `4` contributes 0.} \\
\texttt{        {-} In `{[}2 4 1 5{]}`, `4` contributes 0.} \\
\texttt{        {-} Then `{[}1 2 3 4{]}` has `1, 2, 3` in. `{[}2 4 1 5{]}` has `1, 2` in? No, `4` out. `1, 2, 5` in?} \\
\texttt{        {-} If `4` is out, code is subset of `1, 2, 3, 5, 6`.} \\
\texttt{        {-} `{[}1 2 3 4{]}` (3 pegs) {-}> `1, 2, 3` in.} \\
\texttt{        {-} `{[}2 4 1 5{]}` (2 pegs) {-}> `1, 2` in (since `4` out). `5` out?} \\
\texttt{        {-} If `1, 2, 3` in, and `5` out, then `6` must be in.} \\
\texttt{        {-} Code: `1, 2, 3, 6`.} \\
\texttt{        {-} Check `{[}1 2 3 4{]}`: `1, 2, 3` in. `4` out. 3 pegs. (1B, 2W).} \\
\texttt{        {-} Check `{[}2 4 1 5{]}`: `2, 1` in. `4, 5` out. 2 pegs. (1B, 1W).} \\
\texttt{        {-} This fits perfectly!} \\
\texttt{        {-} So the code is a permutation of **`1, 2, 3, 6`**.} \\
\texttt{        {-} `4` and `5` are out.} \\
\\ 
\texttt{Now let's find the positions for `1, 2, 3, 6`.} \\
\texttt{{-} From `{[}1 2 3 4{]}` (1B, 2W):} \\
\texttt{    {-} `1` at pos 1, `2` at pos 2, `3` at pos 3.} \\
\texttt{    {-} One is Black, two are White.} \\
\texttt{{-} From `{[}2 4 1 5{]}` (1B, 1W):} \\
\texttt{    {-} `2` at pos 1, `1` at pos 3. (`4, 5` out).} \\
\texttt{    {-} One is Black, one is White.} \\
\texttt{{-} From `{[}3 1 2 4{]}` (0B, 3W):} \\
\texttt{    {-} `3` at pos 1, `1` at pos 2, `2` at pos 3.} \\
\texttt{    {-} All three (`1, 2, 3`) are White.} \\
\texttt{    {-} So:} \\
\texttt{        {-} `3` is not at pos 1.} \\
\texttt{        {-} `1` is not at pos 2.} \\
\texttt{        {-} `2` is not at pos 3.} \\
\texttt{{-} From `{[}1 3 5 6{]}` (0B, 2W):} \\
\texttt{    {-} `1` at pos 1, `3` at pos 2, `6` at pos 3 or 4.} \\
\texttt{    {-} `1` is White (not at pos 1).} \\
\texttt{    {-} `3` is White (not at pos 2).} \\
\texttt{    {-} `6` is White (not at pos 3 or 4).} \\
\texttt{    {-} So `6` is not at pos 3 or 4. `6` must be at pos 1 or 2.} \\
\\ 
\texttt{Let's combine constraints:} \\
\texttt{{-} `1` not at pos 1 (from `{[}1 3 5 6{]}`).} \\
\texttt{{-} `1` not at pos 2 (from `{[}3 1 2 4{]}`).} \\
\texttt{{-} So **`1` is at pos 3 or 4**.} \\
\texttt{{-} `2` not at pos 2 (from `{[}1 2 3 4{]}`? No, we don't know if `2` is B there yet).} \\
\texttt{{-} `2` not at pos 3 (from `{[}3 1 2 4{]}`).} \\
\texttt{{-} `2` not at pos 1? In `{[}2 4 1 5{]}`, `2` is at pos 1. If `2` was Black there, then `1` is White.} \\
\texttt{    {-} If `2` is Black in `{[}2 4 1 5{]}` (pos 1):} \\
\texttt{        {-} Then `1` is White in `{[}2 4 1 5{]}` (pos 3).} \\
\texttt{        {-} In `{[}1 2 3 4{]}`, `2` is at pos 1. If `2` is Black, then `1` and `3` are White.} \\
\texttt{        {-} In `{[}3 1 2 4{]}`, `2` is at pos 3. If `2` is at pos 1, it's White. `1` at pos 2 (White). `3` at pos 1 (White). Total 3W. Matches.} \\
\texttt{        {-} So `2` at pos 1 is a strong candidate.} \\
\texttt{{-} `3` not at pos 1 (from `{[}3 1 2 4{]}`).} \\
\texttt{{-} `3` not at pos 2 (from `{[}1 3 5 6{]}`).} \\
\texttt{{-} `3` not at pos 3? In `{[}1 2 3 4{]}`, `3` is at pos 3. If `3` is Black, then `1, 2` are White.} \\
\texttt{    {-} If `3` is Black (pos 3):} \\
\texttt{        {-} In `{[}3 1 2 4{]}`, `3` is at pos 1 (White). `1` at pos 2 (White). `2` at pos 3 (White). Total 3W. Matches.} \\
\texttt{        {-} In `{[}1 2 3 4{]}`, `3` is B. `1, 2` are W. Matches.} \\
\texttt{        {-} In `{[}2 4 1 5{]}`, `3` is not in guess. `1` at pos 3 (W). `2` at pos 1 (W). Total 2W. But we got 1B, 1W.} \\
\texttt{        {-} So `3` cannot be the Black peg in `{[}1 2 3 4{]}`.} \\
\texttt{{-} So the Black peg in `{[}1 2 3 4{]}` is either `1` or `2`.} \\
\texttt{    {-} If `1` is Black (pos 1):} \\
\texttt{        {-} Contradicts `1` not at pos 1 (from `{[}1 3 5 6{]}`).} \\
\texttt{    {-} So **`2` must be the Black peg in `{[}1 2 3 4{]}`**.} \\
\texttt{    {-} Therefore, **`2` is at pos 2**.} \\
\\ 
\texttt{Current State: `{[}? 2 ? ?{]}`} \\
\texttt{Digits: `1, 3, 6`.} \\
\texttt{Constraints:} \\
\texttt{{-} `1` not at pos 1, not at pos 2. So `1` is at pos 3 or 4.} \\
\texttt{{-} `3` not at pos 1, not at pos 2. So `3` is at pos 3 or 4.} \\
\texttt{{-} `6` not at pos 3, not at pos 4 (from `{[}1 3 5 6{]}`). So `6` is at pos 1.} \\
\texttt{{-} Therefore, **`6` is at pos 1**.} \\
\\ 
\texttt{Current State: `{[}6 2 ? ?{]}`} \\
\texttt{Remaining positions: 3, 4.} \\
\texttt{Remaining digits: `1, 3`.} \\
\texttt{{-} `1` is at pos 3 or 4.} \\
\texttt{{-} `3` is at pos 3 or 4.} \\
\texttt{{-} Check `{[}2 4 1 5{]}` (1B, 1W):} \\
\texttt{    {-} `2` is at pos 1. Code has `2` at pos 2. So `2` is White.} \\
\texttt{    {-} We need 1 Black. So either `1` or `5` (out) or `4` (out) must be Black? No, `4, 5` are out.} \\
\texttt{    {-} So `1` must be the Black peg in `{[}2 4 1 5{]}`.} \\
\texttt{    {-} In `{[}2 4 1 5{]}`, `1` is at pos 3.} \\
\texttt{    {-} So **`1` is at pos 3**.} \\
\texttt{{-} Therefore, **`3` is at pos 4**.} \\
\\ 
\texttt{Proposed Code: **`6 2 1 3`**.} \\
\\ 
\texttt{Let's verify:} \\
\texttt{{-} `{[}1 2 3 4{]}`: `1`(W), `2`(B), `3`(W), `4`(out). {-}> 1B, 2W. (Matches)} \\
\texttt{{-} `{[}1 2 5 6{]}`: `1`(W), `2`(W), `5`(out), `6`(W). {-}> 0B, 3W.} \\
\texttt{    {-} Wait, my previous feedback for `{[}1 2 5 6{]}` was **0B, 2W**.} \\
\texttt{    {-} My proposed code `6 2 1 3` gives 3W for `{[}1 2 5 6{]}`.} \\
\texttt{    {-} Contradiction.} \\
\\ 
\texttt{Re{-}evaluate `{[}1 2 5 6{]}` feedback: 0B, 2W.} \\
\texttt{{-} Digits in guess: `1, 2, 5, 6`.} \\
\texttt{{-} Digits in code: `1, 2, 3, 6` (based on `4, 5` out).} \\
\texttt{{-} Intersection: `1, 2, 6`. (3 digits).} \\
\texttt{{-} If 3 digits are in the code, we should get 3 pegs (since 0B).} \\
\texttt{{-} But we got 2W.} \\
\texttt{{-} This means one of `1, 2, 6` is **NOT** in the code.} \\
\texttt{{-} Which one?} \\
\texttt{    {-} If `1` is out: Code `2, 3, 6, 4`? No, `4` out. `2, 3, 6, ?`.} \\
\texttt{    {-} If `2` is out: Code `1, 3, 6, 4`? No, `4` out.} \\
\texttt{    {-} If `6` is out: Code `1, 2, 3, 4`? No, `4` out.} \\
\texttt{    {-} Maybe `4` IS in the code?} \\
\texttt{    {-} Let's re{-}examine `{[}1 2 3 4{]}` (1B, 2W).} \\
\texttt{        {-} If `4` is in, and `1, 2, 3` are in, that's 4 digits.} \\
\texttt{        {-} Then `{[}1 2 3 4{]}` has 4 correct digits. But we got 3 pegs.} \\
\texttt{        {-} So one of `1, 2, 3, 4` is out.} \\
\texttt{    {-} If `4` is in, then one of `1, 2, 3` is out.} \\
\texttt{    {-} If `1` is out: Code `2, 3, 4, 6`.} \\
\texttt{        {-} `{[}1 2 3 4{]}`: `2, 3, 4` in. 3 pegs. (1B, 2W).} \\
\texttt{        {-} `{[}1 2 5 6{]}`: `2, 6` in. `1, 5` out. 2 pegs. (0B, 2W). **Matches!**} \\
\texttt{        {-} `{[}3 1 2 4{]}`: `3, 2, 4` in. 3 pegs. (0B, 3W). **Matches!**} \\
\texttt{        {-} `{[}1 3 5 6{]}`: `3, 6` in. `1, 5` out. 2 pegs. (0B, 2W). **Matches!**} \\
\texttt{        {-} `{[}2 4 1 5{]}`: `2, 4` in. `1, 5` out. 2 pegs. (1B, 1W).} \\
\texttt{            {-} We need 1B, 1W.} \\
\texttt{            {-} Digits `2, 4` are in.} \\
\texttt{            {-} If `2` is Black (pos 1 in guess? No, `2` is at pos 1 in guess).} \\
\texttt{            {-} If `2` is at pos 1 in code:} \\
\texttt{                {-} `{[}2 4 1 5{]}`: `2` (B). `4` (W or B). `1` (out). `5` (out).} \\
\texttt{                {-} If `4` is W, we get 1B, 1W.} \\
\texttt{                {-} Check `{[}1 2 3 4{]}`: `2` at pos 2. If `2` is at pos 1, it's W. `3` at pos 3. `4` at pos 4.} \\
\texttt{                {-} If `2` is W, `3` is W, `4` is B (pos 4).} \\
\texttt{                {-} Then `{[}1 2 3 4{]}` is 1B (`4`), 2W (`2, 3`). Matches.} \\
\texttt{                {-} Check `{[}3 1 2 4{]}`: `3` at pos 1 (W). `1` out. `2` at pos 3 (W). `4` at pos 4 (B).} \\
\texttt{                {-} Result: 1B, 2W. But we got 0B, 3W.} \\
\texttt{                {-} So `4` cannot be Black} \\
            }
        }
    }
    & & \\ \\

    \theutterance \stepcounter{utterance}  
    & & & \multicolumn{4}{p{0.6\linewidth}}{
        \cellcolor[rgb]{0.9,0.9,0.9}{
            \makecell[{{p{\linewidth}}}]{
                \texttt{\tiny{[P1$\langle$GM]}}
                \texttt{{[}GAME{]} You attempted an invalid move. Reason: You have already guessed {[}3, 1, 2, 4{]}. Please try a different guess. Please resubmit a valid move and remember to follow the game rules to avoid penalties.} \\
            }
        }
    }
     \\ \\

    \theutterance \stepcounter{utterance}  
    & \multicolumn{4}{p{0.6\linewidth}}{
        \cellcolor[rgb]{0.561,0.737,0.561}{
            \makecell[{{p{\linewidth}}}]{
                \texttt{\tiny{[P1$\rangle$GM]}}
                \texttt{{[}4 6 2 1{]}} \\
            }
        }
    }
    & & \\ \\

    \theutterance \stepcounter{utterance}  
    & & & \multicolumn{4}{p{0.6\linewidth}}{
        \cellcolor[rgb]{0.9,0.9,0.9}{
            \makecell[{{p{\linewidth}}}]{
                \texttt{\tiny{[P1$\langle$GM]}}
                \texttt{{[}GAME{]} Submitted {[}4 6 2 1{]}. Feedback: 0 black peg(s), 3 white peg(s).} \\
            }
        }
    }
     \\ \\

    \theutterance \stepcounter{utterance}  
    & \multicolumn{4}{p{0.6\linewidth}}{
        \cellcolor[rgb]{0.561,0.737,0.561}{
            \makecell[{{p{\linewidth}}}]{
                \texttt{\tiny{[P1$\rangle$GM]}}
                \texttt{Based on the feedback from `{[}4 6 2 1{]}` (0 Black, 3 White):} \\
\texttt{{-} **0 Black**: None of the digits are in the correct position.} \\
\texttt{    {-} `4` is not at pos 1.} \\
\texttt{    {-} `6` is not at pos 2.} \\
\texttt{    {-} `2` is not at pos 3.} \\
\texttt{    {-} `1` is not at pos 4.} \\
\texttt{{-} **3 White**: Three digits from my guess are in the code.} \\
\texttt{{-} **Deduction on Digits**:} \\
\texttt{    {-} My guess contains `1, 2, 4, 6`.} \\
\texttt{    {-} We have 3 correct digits.} \\
\texttt{    {-} Let's re{-}evaluate the set of digits using all feedback.} \\
\texttt{    {-} Guess 1: `{[}1 2 3 4{]}` {-}> 1B, 2W (3 correct).} \\
\texttt{    {-} Guess 2: `{[}1 2 5 6{]}` {-}> 0B, 2W (2 correct).} \\
\texttt{    {-} Guess 3: `{[}3 1 2 4{]}` {-}> 0B, 3W (3 correct).} \\
\texttt{    {-} Guess 4: `{[}1 3 5 6{]}` {-}> 0B, 2W (2 correct).} \\
\texttt{    {-} Guess 5: `{[}2 4 1 5{]}` {-}> 1B, 1W (2 correct).} \\
\texttt{    {-} Guess 6: `{[}4 6 2 1{]}` {-}> 0B, 3W (3 correct).} \\
\\ 
\texttt{    {-} **Intersection Analysis**:} \\
\texttt{        {-} Compare Guess 2 (`1, 2, 5, 6`, 2 correct) and Guess 4 (`1, 3, 5, 6`, 2 correct).} \\
\texttt{            {-} Common: `1, 5, 6`.} \\
\texttt{            {-} If `1` is in the code, then `5` and `6` must be out (otherwise Guess 4 would have 3 correct).} \\
\texttt{            {-} If `1` is out, then `5` and `6` must be in (to get 2 correct in both).} \\
\texttt{            {-} **Hypothesis A**: `1` is in, `5, 6` out.} \\
\texttt{                {-} Then Guess 2 (`1, 2, 5, 6`) has `1` and `2` (since `5, 6` out). So `2` is in.} \\
\texttt{                {-} Guess 4 (`1, 3, 5, 6`) has `1` and `3` (since `5, 6` out). So `3` is in.} \\
\texttt{                {-} So far: `1, 2, 3` are in.} \\
\texttt{                {-} Guess 1 (`1, 2, 3, 4`) has `1, 2, 3` in. So `4` is out.} \\
\texttt{                {-} Code digits: `1, 2, 3` and one more. The only remaining digit is `4` (out), `5` (out), `6` (out). This is impossible. We need 4 digits.} \\
\texttt{                {-} **Conclusion**: Hypothesis A is false. `1` is **NOT** in the code.} \\
\texttt{            {-} **Hypothesis B**: `1` is out, `5, 6` are in.} \\
\texttt{                {-} Guess 2 (`1, 2, 5, 6`) has `5, 6` in. So `2` is out.} \\
\texttt{                {-} Guess 4 (`1, 3, 5, 6`) has `5, 6` in. So `3` is out.} \\
\texttt{                {-} So far: `5, 6` are in. `1, 2, 3` are out.} \\
\texttt{                {-} Remaining digits: `4`. We need 2 more digits to make 4. But we only have `4` left? No, digits are 1{-}6.} \\
\texttt{                {-} If `1, 2, 3` are out, the code must be `4, 5, 6` and... wait, we need 4 digits.} \\
\texttt{                {-} This implies my deduction about "out" is too strict. Let's re{-}evaluate the counts.} \\
\texttt{                {-} Maybe `1` is in, but `5` or `6` is also in?} \\
\texttt{                {-} Let's look at Guess 6 (`4 6 2 1`) {-}> 3W.} \\
\texttt{                    {-} Digits: `4, 6, 2, 1`.} \\
\texttt{                    {-} If `1` is out (from Hyp B), then `4, 6, 2` are in. That's 3 digits. Matches 3W.} \\
\texttt{                    {-} So `1` is out. `4, 6, 2` are in.} \\
\texttt{                    {-} But Hyp B said `2` is out. Contradiction.} \\
\texttt{                {-} Let's re{-}evaluate Hyp A: `1` is in, `5, 6` out.} \\
\texttt{                    {-} Guess 6 (`4 6 2 1`) has `1` in. `6` out. So `4, 2` must be in to get 3W.} \\
\texttt{                    {-} So `1, 2, 4` are in. `6` is out.} \\
\texttt{                    {-} Check Guess 2 (`1 2 5 6`): `1, 2` in. `5, 6` out. Total 2 correct. Matches 2W.} \\
\texttt{                    {-} Check Guess 4 (`1 3 5 6`): `1` in. `5, 6` out. So `3` must be in to get 2W.} \\
\texttt{                    {-} So `1, 3` in.} \\
\texttt{                    {-} So far: `1, 2, 3, 4` are in. `5, 6` out.} \\
\texttt{                    {-} Check Guess 1 (`1 2 3 4`): `1, 2, 3, 4` in. Total 4 correct. But we got 3 pegs (1B, 2W).} \\
\texttt{                    {-} This implies one of `1, 2, 3, 4` is NOT in the code.} \\
\texttt{                    {-} But we just deduced `1, 2, 3, 4` are all in.} \\
\texttt{                    {-} Where is the error?} \\
\texttt{                    {-} Error in "If `1` is in, then `5` and `6` must be out".} \\
\texttt{                    {-} In Guess 4 (`1 3 5 6`), if `1` is in, and we have 2 pegs.} \\
\texttt{                    {-} If `5` and `6` are also in, we have 3 correct digits (`1, 3, 5` or `1, 3, 6`).} \\
\texttt{                    {-} If we have 3 correct digits, we should get 3 pegs (since 0B). But we got 2W.} \\
\texttt{                    {-} So indeed, if `1` is in, we cannot have both `5` and `6` in.} \\
\texttt{                    {-} Can we have `1` in, `5` in, `6` out?} \\
\texttt{                        {-} Guess 2 (`1 2 5 6`): `1, 5` in. `2, 6`? If `2` out, then 2 correct. Matches.} \\
\texttt{                        {-} Guess 4 (`1 3 5 6`): `1, 5` in. `3, 6`? If `3` out, then 2 correct. Matches.} \\
\texttt{                        {-} So `1, 5` in. `2, 3, 6` out?} \\
\texttt{                        {-} Check Guess 1 (`1 2 3 4`): `1` in. `2, 3` out. So `4` must be in to get 3 pegs.} \\
\texttt{                        {-} So `1, 4, 5` in. `2, 3, 6` out.} \\
\texttt{                        {-} Check Guess 6 (`4 6 2 1`): `4, 1` in. `6, 2` out. Total 2 correct. But we got 3W.} \\
\texttt{                        {-} Contradiction.} \\
\texttt{                    {-} Can we have `1` in, `6` in, `5` out?} \\
\texttt{                        {-} Guess 2 (`1 2 5 6`): `1, 6` in. `2, 5`? If `2` out, 2 correct. Matches.} \\
\texttt{                        {-} Guess 4 (`1 3 5 6`): `1, 6` in. `3, 5`? If `3` out, 2 correct. Matches.} \\
\texttt{                        {-} So `1, 6` in. `2, 3, 5` out?} \\
\texttt{                        {-} Check Guess 1 (`1 2 3 4`): `1` in. `2, 3` out. So `4` must be in.} \\
\texttt{                        {-} So `1, 4, 6` in. `2, 3, 5` out.} \\
\texttt{                        {-} Check Guess 6 (`4 6 2 1`): `4, 6, 1` in. `2` out. Total 3 correct. Matches 3W.} \\
\texttt{                        {-} This looks promising!} \\
\texttt{                        {-} Code digits: **`1, 4, 6`** and one more.} \\
\texttt{                        {-} The missing digit must be from `\{2, 3, 5\}`.} \\
\texttt{                        {-} We assumed `2, 3, 5` are out to satisfy Guess 2 and 4.} \\
\texttt{                        {-} But we need 4 digits. So one of `2, 3, 5` MUST be in.} \\
\texttt{                        {-} If `2` is in:} \\
\texttt{                            {-} Guess 2 (`1 2 5 6`): `1, 2, 6` in. 3 correct. Should be 3W. But we got 2W.} \\
\texttt{                            {-} So `2` cannot be in.} \\
\texttt{                        {-} If `3` is in:} \\
\texttt{                            {-} Guess 4 (`1 3 5 6`): `1, 3, 6` in. 3 correct. Should be 3W. But we got 2W.} \\
\texttt{                            {-} So `3` cannot be in.} \\
\texttt{                        {-} If `5` is in:} \\
\texttt{                            {-} Guess 2 (`1 2 5 6`): `1, 5, 6` in. 3 correct. Should be 3W. But we got 2W.} \\
\texttt{                            {-} So `5` cannot be in.} \\
\texttt{                        {-} **Contradiction again.**} \\
\\ 
\texttt{    {-} Let's rethink the "3 correct digits" logic.} \\
\texttt{    {-} Maybe `1` is NOT in the code?} \\
\texttt{    {-} If `1` is out:} \\
\texttt{        {-} Guess 2 (`1 2 5 6`): `2, 5, 6`? We need 2 correct.} \\
\texttt{        {-} Guess 4 (`1 3 5 6`): `3, 5, 6`? We need 2 correct.} \\
\texttt{        {-} Guess 6 (`4 6 2 1`): `4, 6, 2`? We need 3 correct.} \\
\texttt{        {-} From Guess 6: `4, 6, 2` are in. (Since `1` is out).} \\
\texttt{        {-} So `2, 4, 6` are in.} \\
\texttt{        {-} Check Guess 2 (`1 2 5 6`): `2, 6` in. `1, 5` out. Total 2 correct. Matches 2W.} \\
\texttt{        {-} Check Guess 4 (`1 3 5 6`): `6` in. `1, 5` out. So `3` must be in to get 2 correct.} \\
\texttt{        {-} So `3` is in.} \\
\texttt{        {-} So far: `2, 3, 4, 6` are in. `1, 5` are out.} \\
\texttt{        {-} Check Guess 1 (`1 2 3 4`): `2, 3, 4` in. `1` out. Total 3 correct. Matches 3 pegs (1B, 2W).} \\
\texttt{        {-} Check Guess 5 (`2 4 1 5`): `2, 4` in. `1, 5` out. Total 2 correct. Matches 2 pegs (1B, 1W).} \\
\texttt{        {-} Check Guess 3 (`3 1 2 4`): `3, 2, 4` in. `1` out. Total 3 correct. Matches 3 pegs (0B, 3W).} \\
\texttt{        {-} **This set works perfectly!**} \\
\texttt{        {-} **Code Digits**: **`2, 3, 4, 6`**.} \\
\texttt{        {-} **Out**: `1, 5`.} \\
\\ 
\texttt{    {-} Now determine positions for `2, 3, 4, 6`.} \\
\texttt{    {-} Constraints:} \\
\texttt{        {-} Guess 1 (`1 2 3 4`): 1B, 2W.} \\
\texttt{            {-} `2` at pos 2, `3` at pos 3, `4` at pos 4.} \\
\texttt{            {-} One is Black, two are White.} \\
\texttt{        {-} Guess 3 (`3 1 2 4`): 0B, 3W.} \\
\texttt{            {-} `3` at pos 1, `2` at pos 3, `4` at pos 4.} \\
\texttt{            {-} All are White.} \\
\texttt{            {-} So: `3` not at 1, `2` not at 3, `4` not at 4.} \\
\texttt{        {-} Guess 5 (`2 4 1 5`): 1B, 1W.} \\
\texttt{            {-} `2` at pos 1, `4` at pos 2. (`1, 5` out).} \\
\texttt{            {-} One is Black, one is White.} \\
\texttt{        {-} Guess 6 (`4 6 2 1`): 0B, 3W.} \\
\texttt{            {-} `4` at pos 1, `6` at pos 2, `2` at pos 3.} \\
\texttt{            {-} All are White.} \\
\texttt{            {-} So: `4` not at 1, `6` not at 2, `2` not at 3.} \\
\\ 
\texttt{    {-} Combine constraints:} \\
\texttt{        {-} `2` not at 3 (from G3 and G6).} \\
\texttt{        {-} `4` not at 4 (from G3) and `4` not at 1 (from G6).} \\
\texttt{        {-} `3` not at 1 (from G3).} \\
\texttt{        {-} `6` not at 2 (from G6).} \\
\\ 
\texttt{    {-} From Guess 5 (`2 4 1 5`):} \\
\texttt{        {-} `2` at pos 1, `4` at pos 2.} \\
\texttt{        {-} One is Black.} \\
\texttt{        {-} If `2` is Black (pos 1):} \\
\texttt{            {-} Then `4` is White (pos 2).} \\
\texttt{            {-} So `2` is at pos 1.} \\
\texttt{            {-} `4` is not at pos 2.} \\
\texttt{        {-} If `4` is Black (pos 2):} \\
\texttt{            {-} Then `2` is White (pos 1).} \\
\texttt{            {-} So `4` is at pos 2.} \\
\texttt{            {-} `2` is not at pos 1.} \\
\\ 
\texttt{    {-} Let's test **Case A: `2` is at pos 1**.} \\
\texttt{        {-} Code: `{[}2 ? ? ?{]}`.} \\
\texttt{        {-} `4` is not at 1, 2, 4. So `4` must be at pos 3.} \\
\texttt{        {-} Code: `{[}2 ? 4 ?{]}`.} \\
\texttt{        {-} `3` not at 1. `3` is at 2 or 4.} \\
\texttt{        {-} `6` not at 2. `6` is at 4 (since 1, 3 taken).} \\
\texttt{        {-} So `6` is at pos 4.} \\
\texttt{        {-} Code: `{[}2 ? 4 6{]}`.} \\
\texttt{        {-} `3` must be at pos 2.} \\
\texttt{        {-} Proposed Code: **`2 3 4 6`**.} \\
\texttt{        {-} Verify Guess 1 (`1 2 3 4`):} \\
\texttt{            {-} `2` at pos 2 (Code has `2` at 1) {-}> W.} \\
\texttt{            {-} `3` at pos 3 (Code has `3` at 2) {-}> W.} \\
\texttt{            {-} `4` at pos 4 (Code has `4` at 3) {-}> W.} \\
\texttt{            {-} Result: 0B, 3W.} \\
\texttt{            {-} But we need 1B, 2W.} \\
\texttt{            {-} **Contradiction**.} \\
\\ 
\texttt{    {-} Let's test **Case B: `4` is at pos 2**.} \\
\texttt{        {-} Code: `{[}? 4 ? ?{]}`.} \\
\texttt{        {-} `2` is not at 1 (from G5 logic: if `4` is B, `2` is W).} \\
\texttt{        {-} `2` not at 3 (from G3, G6).} \\
\texttt{        {-} `2` not at 2 (occupied by `4`).} \\
\texttt{        {-} So `2` must be at pos 4.} \\
\texttt{        {-} Code: `{[}? 4 ? 2{]}`.} \\
\texttt{        {-} `4` not at 1 (from G6).} \\
\texttt{        {-} So `4` is at pos 2. (Consistent).} \\
\texttt{        {-} `3` not at 1 (from G3).} \\
\texttt{        {-} `3` not at 2 (occupied).} \\
\texttt{        {-} `3` not at 4 (occupied).} \\
\texttt{        {-} So `3` must be at pos 3.} \\
\texttt{        {-} Code: `{[}? 4 3 2{]}`.} \\
\texttt{        {-} `6` must be at pos 1.} \\
\texttt{        {-} Proposed Code: **`6 4 3 2`**.} \\
\texttt{        {-} Verify Guess 1 (`1 2 3 4`):} \\
\texttt{            {-} `2` at pos 2 (Code `2` at 4) {-}> W.} \\
\texttt{            {-} `3` at pos 3 (Code `3` at 3) {-}> **B**.} \\
\texttt{            {-} `4` at pos 4 (Code `4` at 2) {-}> W.} \\
\texttt{            {-} Result: 1B, 2W. **Matches!**} \\
\texttt{        {-} Verify Guess 3 (`3 1 2 4`):} \\
\texttt{            {-} `3` at pos 1 (Code `3` at 3) {-}> W.} \\
\texttt{            {-} `2` at pos 3 (Code `2` at 4) {-}> W.} \\
\texttt{            {-} `4` at pos 4 (Code `4` at 2) {-}> W.} \\
\texttt{            {-} Result: 0B, 3W. **Matches!**} \\
\texttt{        {-} Verify Guess 5 (`2 4 1 5`):} \\
\texttt{            {-} `2` at pos 1 (Code `2` at 4) {-}> W.} \\
\texttt{            {-} `4` at pos 2 (Code `4` at 2) {-}> **B**.} \\
\texttt{            {-} Result: 1B, 1W. **Matches!**} \\
\texttt{        {-} Verify Guess 6 (`4 6 2 1`):} \\
\texttt{            {-} `4` at pos 1 (Code `4` at 2) {-}> W.} \\
\texttt{            {-} `6` at pos 2 (Code `6` at 1) {-}> W.} \\
\texttt{            {-} `2` at pos 3 (Code `2` at 4) {-}> W.} \\
\texttt{            {-} Result: 0B, 3W. **Matches!**} \\
\texttt{        {-} Verify Guess 2 (`1 2 5 6`):} \\
\texttt{            {-} `2` at pos 2 (Code `2` at 4) {-}> W.} \\
\texttt{            {-} `6` at pos 4 (Code `6` at 1) {-}> W.} \\
\texttt{            {-} Result: 0B, 2W. **Matches!**} \\
\texttt{        {-} Verify Guess 4 (`1 3 5 6`):} \\
\texttt{            {-} `3` at pos 2 (Code `3` at 3) {-}> W.} \\
\texttt{            {-} `6` at pos 4 (Code `6` at 1) {-}> W.} \\
\texttt{            {-} Result: 0B, 2W. **Matches!**} \\
\\ 
\texttt{    {-} The code is **`6 4 3 2`**.} \\
\\ 
\texttt{{[}6 4 3 2{]}} \\
            }
        }
    }
    & & \\ \\

    \theutterance \stepcounter{utterance}  
    & & & \multicolumn{2}{p{0.3\linewidth}}{
        \cellcolor[rgb]{0.9,0.9,0.9}{
            \makecell[{{p{\linewidth}}}]{
                \texttt{\tiny{[GM$|$GM]}}
                \texttt{role: Player 0} \\
\texttt{invalid\_move: False} \\
\texttt{turn\_count: 9} \\
\texttt{reason: You have cracked the code, solving 4 out of 4 pegs correctly} \\
\texttt{reward: 1} \\
            }
        }
    }
    & & \\ \\

\end{supertabular}
}

\end{document}
